% !TeX program = pdflatex
% Orchestration Economics — DBA Thesis Draft v14 (business-reader narrative restructure)
% Youssef Hariri, DeeWaanAI, 2026-09-05
\documentclass[12pt,a4paper]{article}

% ── Packages ──────────────────────────────────────────────────────
\usepackage[utf8]{inputenc}
\usepackage[T1]{fontenc}
\usepackage{lmodern}
\usepackage[margin=2.5cm]{geometry}
\usepackage{amsmath,amssymb}
\usepackage{graphicx}
\usepackage{booktabs}
\usepackage{longtable}
\usepackage{array}
\usepackage{multirow}
\usepackage{hyperref}
\usepackage{natbib}
\usepackage{caption}
\usepackage{subcaption}
\usepackage{enumitem}
\usepackage{xcolor}
\usepackage{mdframed}
\usepackage{csquotes}  % proper quotation marks
\usepackage{siunitx}
\usepackage{float}
\usepackage{tabularx}
\usepackage[capitalise,nameinlink]{cleveref}

% ── Custom commands ───────────────────────────────────────────────
\newcommand{\phiBlend}{\varphi}
\newcommand{\BF}{\ensuremath{\mathit{BF}}}
\newcommand{\TEGR}{\ensuremath{\text{T-EGR}}}
\newcommand{\dval}[1]{\ensuremath{d = #1}}
\newcommand{\CI}[1]{\ensuremath{#1}}
\newcommand{\figref}[1]{Figure~\ref{#1}}
\newcommand{\tabref}[1]{Table~\ref{#1}}
\newcommand{\secref}[1]{Section~\ref{#1}}

% ── Executive-summary box ─────────────────────────────────────────
\mdfdefinestyle{execsummary}{
  linecolor=black,
  linewidth=1.5pt,
  backgroundcolor=gray!5,
  innertopmargin=10pt,
  innerbottommargin=10pt,
  innerrightmargin=10pt,
  innerleftmargin=10pt,
  skipabove=\bigskipamount,
  skipbelow=\bigskipamount,
}

% ── Hyperref setup ────────────────────────────────────────────────
\hypersetup{
  colorlinks=true,
  linkcolor=blue!60!black,
  citecolor=green!40!black,
  urlcolor=blue!70!black,
}

% ── Title ─────────────────────────────────────────────────────────
\title{\textbf{Orchestrate or Upgrade? When a Lower-Tier Model Delivers
98\% Quality at 23\% Cost}}

\author{Youssef Hariri\\
\small DeeWaanAI\quad $\cdot$\quad \href{https://osf.io/mjb8x}{\texttt{osf.io/mjb8x}}\quad $\cdot$\quad 2026-08-28}

\date{}

% ==================================================================
\begin{document}
\maketitle

\begin{abstract}
\noindent
Organizations deploying large language models (LLMs) for knowledge work face a
budget choice: pay for a frontier model, or orchestrate cheaper models into a
multi-step pipeline. This thesis answers that choice with three pre-registered
studies that escalate a single question.

\smallskip
\noindent\textbf{Study 1 --- Does structure substitute for tier?} In a
$2 \times 2 \times 2$ factorial across five LLM families and 19 business
problems (1{,}900 blinded generation--judgment runs), adding a three-step team
structure lifts a lower-tier model to \textbf{98.8\% of frontier quality} at
88\% of the cost ($\TEGR = 1.49$), and structure helps cheaper models
significantly more than frontier ones ($d = -0.227$, $p = 0.015$). Yet a
\emph{single} lower-tier call already captures 98\% of frontier quality at just
23\% of the cost ($\TEGR = 6.15$) --- the highest-return configuration tested ---
so the team's added structure buys only a marginal quality gain for about
3.8$\times$ the spend. Frontier teams are the worst dollar (3.5$\times$ the cost
for negligible gain, $\TEGR = 0.27$). A pre-registered replication on 30 fresh problems and 10 families confirms the core economics out-of-sample; three secondary hypotheses --- a marketing-domain exception, a capability-gap correlation, and team-driven variance reduction --- failed their pre-registered gates and are not recommended.

\smallskip
\noindent\textbf{Study 2 --- Why substitution stops short.} A \$0.002
refinement pass recovers up to 64\% of the frontier gap but never achieves
paired superiority. Qualitative coding ($\kappa = 0.824$) locates the ceiling:
orchestration buys \emph{coverage} (more structure, more length) while the
frontier wins on \emph{density} (compressed insight, zero boilerplate) --- a gap
that add-only pipelines cannot close.

\smallskip
\noindent\textbf{Study 3 --- Who should design the pipeline?} When models
auto-design their own workflows, authorship barely matters: self-designed and
frontier-designed pipelines are statistically indistinguishable in four of five
families, and both beat our frozen human-designed instrument. But orchestration
economics are family-specific --- at 0.24--3.5$\times$ solo cost, cheap
auto-designed pipelines beat the frontier-solo baseline for Kimi ($d = +0.85$ on
novelty, the only family to survive multiplicity control) and lose decisively for
Claude ($d = -3.15$). A pre-registered replication on 30 fresh problems finds the
pattern transports only at its extremes (Kimi $d = +1.25$; Claude $d = -1.90$),
so routing must be \emph{measured per deployment, not looked up}.

\smallskip
\noindent\textbf{For business decision-makers:} use a \emph{single} lower-tier
call by default --- the highest return. Add a \emph{team} (structure) only when
the binding gap is \emph{coverage} (breadth, length, completeness) and its
roughly four-fold cost is justified. Reserve the \emph{frontier} for
\emph{density}-critical work --- compressed insight with zero boilerplate --- the
one thing orchestration cannot manufacture. Never combine frontier and team (the
worst dollar in the study). And because these economics are family-specific,
\emph{measure} efficiency ($\TEGR$) on your own models and problems rather than
trusting any lookup table.

\medskip
\noindent\textbf{Keywords:} Large Language Models, Multi-Agent Systems, LLM Orchestration,
AI Cost Optimization, Pipeline Architecture, LLM-as-a-Judge, Resource
Allocation, Pre-registered Experiment, Cost-Effectiveness Analysis, Business
AI Strategy, T-EGR, Model Tier Substitution.
\end{abstract}

\bigskip

\noindent\textit{A practitioner-oriented executive summary with budget-allocation
rules appears at the end of the paper (\cref{sec:exec-summary}).}

% ── List of Abbreviations ─────────────────────────────────────────
\section*{List of Abbreviations}
\small
\begin{longtable}{@{}ll@{}}
\toprule
\textbf{Abbr} & \textbf{Definition} \\
\midrule
$\BF_{10}$        & Bayes factor in favor of $H_1$ over $H_0$ \\
C-S-N             & lower-tier solo, no persona \\
C-S-P             & lower-tier solo, persona \\
C-T-N             & lower-tier team, no persona \\
C-T-P             & lower-tier team, persona \\
F-S-N             & frontier solo, no persona \\
F-S-P             & frontier solo, persona \\
F-T-N             & frontier team, no persona \\
F-T-P             & frontier team, persona \\
EGR               & Efficiency Gain Ratio \\
FDR               & false discovery rate \\
FLOPs             & floating-point operations \\
FSN               & frontier solo, no persona (baseline) \\
ICER              & incremental cost--effectiveness ratio \\
IRR               & inter-rater reliability \\
JZS               & Jeffreys--Zellner--Siow (Bayes factor prior) \\
LLM               & large language model \\
LMM               & linear mixed model \\
LOO               & leave-one-out \\
OBF               & O'Brien--Fleming (spending function) \\
SESOI             & smallest effect size of interest \\
T-EGR             & Tiered Efficiency Gain Ratio \\
TOST              & two one-sided tests (equivalence) \\
\bottomrule
\end{longtable}
\normalsize

% ==================================================================
% 1  INTRODUCTION
% ==================================================================
\section{Introduction}\label{sec:intro}
Recent work has moved from single-model inference to the orchestration of
multiple language-model calls: refinement pipelines chain builder, critic, and
refiner steps into structured workflows \citep{WeiEtAl2022, ShinnEtAl2023,
YaoEtAl2023}; agent teams assign social roles and have models debate toward an
answer \citep{LiangEtAl2024}; and panels of judges aggregate verdicts across
models \citep{ZhengEtAl2023, VergaEtAl2024}. Across these designs runs an
implicit economic claim --- that \emph{structure can substitute for model tier}.
A practitioner who cannot afford a frontier model, the argument runs, can reach
comparable output quality by composing cheaper models into an orchestrated
system. If true, the claim reshapes how organizations price, budget, and
architect generative-AI products. If false, orchestration is at best a marginal
convenience and at worst a cost multiplier that never reaches frontier quality.

Prior empirical work on generative AI is encouraging but silent on this
specific claim. Controlled productivity studies show that a single generative
assistant raises output quality and reduces task time \citep{NoyZhang2023,
BrynjolfssonLiRaymond2023}, yet they treat the model as a given, single-agent
capability. A field experiment with consultants documented a ``jagged
technological frontier'': large gains inside the model's capability boundary,
large losses outside it \citep{DellAcquaEtAl2023}. Neither study varies the
structure around the model. Multi-agent and ensemble work shows that sampling,
debate, and tree search can improve accuracy \citep{WangEtAl2023,
LiangEtAl2024, YaoEtAl2023}, but the gains are reported without cost-matched
comparison to a stronger single model. A nascent theory of test-time compute
predicts that, for fixed problem difficulty, more inference compute on a
smaller model can outperform less inference compute on a larger one
\citep{SnellEtAl2024}, but its predictions are in units of FLOPs, not dollars,
and only on verifiable tasks.

Four gaps follow. First, no factorial, cost-matched comparison of orchestrated
lower-tier models versus frontier solo exists on open-ended tasks. Second, no
mechanistic account explains \emph{why} orchestration succeeds or fails. Third, no economic
metric is robust to the price-assumption choices that inevitably shape per-run
costing. Fourth, orchestration studies almost invariably freeze a
\emph{human-designed} pipeline: whether models can design their own workflows,
and whether the resulting cost--quality economics transport across model
families and fresh problems, is largely untested --- leaving practitioners
unable to tell whether a winning configuration can be copied or must be
re-measured.

This paper addresses all four. We run a pre-registered $2 \times 2 \times 2$
(tier $\times$ structure $\times$ persona) factorial experiment over five model
families. The third factor --- a fixed cultural persona applied to each pipeline role,
carried from our prior work on persona economics \citep{Hariri2026b} --- is a
secondary, boundary question rather than part of the core tier--structure claim;
it produces a large negative effect, with personas raising cost while degrading
quality (\cref{sec:persona-penalty}). Nineteen confirmatory business problems are drawn from a 30-problem
domain-stratified universe, with 11 problems reserved for feasibility gating per
the pre-registered protocol --- a separation that keeps feasibility probes from
anchoring confirmatory inference \citep{AlbersLakens2018}. The experiment
produces 1{,}900 generation--judgment runs, measured by a blinded three-judge
panel from families disjoint from every generator. Cost is stored as tokens and
re-priced through a frozen price vector under three blended input/output ratios,
so every reported dollar is reproducible and re-priceable. A second study applies
a cheaper, add-only refinement pass to the orchestrated outputs and codes the
qualitative differences between original, refined, and frontier answers, testing
the mechanistic prediction that orchestration buys coverage while the frontier
wins on density. A third study unfreezes the pipeline itself: across the same
five families, models auto-design their own workflows, which we compare against
frontier-designed and our frozen human-designed instruments, and we test whether
the resulting cost--quality economics transport to 30 fresh problems in a
pre-registered replication.

Six contributions follow:
\begin{enumerate}[leftmargin=*,nosep]
  \item \textbf{The first cost-matched, pre-registered factorial} testing whether
        orchestration structure substitutes for model tier on open-ended business
        tasks --- with negative superiority results reported as primary findings,
        not buried.
  \item \textbf{A new economic metric, T-EGR} (quality per dollar, relative to the
        frontier baseline) and a cost--quality Pareto frontier shown to be robust
        to price assumptions.
  \item \textbf{A mechanistic explanation} of why substitution stops short:
        orchestration buys \emph{coverage} while the frontier wins on
        \emph{density} --- a gap a low-cost refinement pass narrows but never
        closes.
  \item \textbf{A pipeline-authorship finding}: who designs the workflow barely
        matters, but the economics are family-specific, so orchestration must be
        \emph{measured per deployment, not looked up}.
  \item \textbf{Practitioner contingency rules} mapping each budget tier and task
        type to the pipeline architecture with the highest expected return.
  \item \textbf{External validity --- and its boundary}: across two confirmatory
        replications on fresh problems and a widened family universe, the core
        cost--quality economics generalize, while the contingent nuances do not,
        identifying which orchestration findings practitioners can rely on.
\end{enumerate}

The paper proceeds as follows. \cref{sec:lit} reviews related literature.
\cref{sec:method} describes the pre-registered design, apparatus, and analysis,
including an AI-assisted research disclosure (\cref{sec:ai-disclosure}).
\cref{sec:results} presents the three studies as a single buyer's decision
journey: Study~1 asks whether structure substitutes for model tier, Study~2
asks why substitution stops short of the frontier, and Study~3 asks who should
design the pipeline. \cref{sec:discussion} discusses theoretical and practical
implications; \cref{sec:managerial} translates the findings into budget rules
and an operating model for practitioners. \cref{sec:limitations} acknowledges
limitations. \cref{sec:future} outlines future directions.
\cref{sec:conclusion} concludes.

% ==================================================================
% 2  LITERATURE REVIEW
% ==================================================================
\section{Literature Review}\label{sec:lit}

This review synthesizes four streams of research: (i)~measured productivity
effects of generative AI, (ii)~multi-agent and ensemble generation methods,
(iii)~LLM-as-a-judge evaluation methodology, and (iv)~open-science methods for
economic measurement. We then integrate these streams with our own prior work on
personas and efficiency metrics to articulate the research gap that motivates
this study.

\subsection{Productivity effects of generative AI}

\subsubsection{Measured gains in knowledge work}

Three controlled studies establish that generative AI raises knowledge-work
productivity, but they also reveal that the gains are conditional and bounded.
\citet{NoyZhang2023} randomized 444 mid-level professionals to ChatGPT-assisted
and control arms on professional writing tasks: assistance reduced completion
time by 0.8 standard deviations and raised output quality by 0.4 standard
deviations, with the largest quality gains among initially lower-performing
workers --- a compression of the performance distribution.
\citet{BrynjolfssonLiRaymond2023} observed 5{,}179 customer-support agents
using a generative assistant in the field and found a 14\% average productivity
increase, concentrated among novice and low-skilled workers (+34\%), consistent
with the model diffusing behaviors of high-ability agents across the workforce.

A critical limitation unites these studies: both treat the model as a
\emph{given, single-agent capability}. The unit of analysis is one model call,
not an orchestrated system. The professional-writing and customer-support tasks
they study are also relatively closed-form: the desired output is
well-specified (a support response, a draft document), and quality can be
partially verified against ground truth. Open-ended strategic tasks --- the
domain of this paper --- have no verifiable correct answer, which changes both
the nature of the ``frontier'' and the applicability of the productivity gains
documented above.

\subsubsection{The jagged frontier and its boundary}

The boundary of these gains is the central caution. \citet{DellAcquaEtAl2023}
ran a field experiment with 758 consultants on realistic professional problems
and documented a ``jagged technological frontier'': inside the model's
capability boundary, consultants completed 12.2\% more tasks, 25.1\% faster,
with 40\% higher quality; outside it, consultants using the model were 19
percentage points less likely to produce correct solutions than controls. Two
productive behavioral strategies emerged (``centaurs'' dividing work at the
boundary, ``cyborgs'' interleaving with the model), but neither changed the
location of the boundary itself.

The jagged frontier concept is especially consequential for open-ended tasks.
In closed-form domains (arithmetic, code, support scripts), the boundary is
sharp: a solution is either correct or not. In open-ended business strategy,
the boundary is diffuse --- there is no single correct answer, only better or
worse responses. This diffuseness creates a fundamental ambiguity: when a cheaper
model produces a plausible but suboptimal strategy, the user may not know they
are ``outside the frontier'' until the solution fails in deployment. The open
question this paper addresses is whether \emph{orchestration structure} can
shift the effective capability frontier for these open-ended tasks --- and at
what cost.

\subsection{Multi-agent systems, reasoning, and inference-time compute}

\subsubsection{From chain-of-thought to structured reasoning}

The modern reasoning literature begins with the observation that language models
solve complex problems better when they generate intermediate reasoning steps.
\citet{WeiEtAl2022} introduced chain-of-thought (CoT) prompting: appending
exemplars that show step-by-step reasoning before the final answer elicited
comparable reasoning in the model, dramatically improving performance on
arithmetic, commonsense, and symbolic reasoning tasks. The CoT paradigm
reframed language-model inference from single-shot generation to
\emph{structured multi-step deliberation} --- a conceptual shift that enabled
the orchestration methods reviewed below.

CoT's direct descendant in the accuracy--cost tradeoff space is
self-consistency. \citet{WangEtAl2023} showed that sampling a diversity of
reasoning paths from a single model and taking a majority vote over final
answers improves arithmetic, commonsense, and symbolic reasoning well beyond
CoT decoding alone. The result establishes the general principle that parallel
sampled calls can be converted into accuracy. Its direct corollary --- that
accuracy is bought at a multiplicative cost in calls and tokens --- is rarely
priced against the alternative of a single stronger model.

\subsubsection{Interactive deliberation: debate, tree search, and agent loops}

Moving beyond independent sampling, recent work has explored \emph{interactive}
structures in which model instances influence each other. \citet{LiangEtAl2024}
showed that multiple LLM instances proposing answers, critiquing each other's
responses, and refining their own over debate rounds improves factuality and
reasoning on arithmetic and knowledge benchmarks, encourages divergent
thinking, and reliably converges to shared answers. \citet{YaoEtAl2023} extended
the reasoning structure from linear chains to branching trees: Tree of Thoughts
treats each reasoning step as a node in a search graph, allowing the model to
explore multiple reasoning paths, backtrack from dead ends, and select the most
promising branch --- deliberate problem solving rather than single-path
generation.

These works demonstrate that structure (sampling, aggregation, critique loops,
tree search) adds quality beyond a single call. What they do not provide is a
cost-matched factorial comparison: prior multi-agent studies report quality
gains without measuring dollars, and without comparing orchestrated lower-tier models
against an un-orchestrated frontier model on identical problems. The
substitution claim --- that structure can stand in for model tier --- is
therefore assumed, not tested.

\subsubsection{The test-time compute hypothesis}

A nascent theoretical framework reframes the model-tier-versus-structure
tradeoff as an allocation problem. \citet{SnellEtAl2024} demonstrated that
scaling inference-time computation (more tokens, verifier-based search,
iterative refinement) can be more effective than scaling model parameters for
improving performance on reasoning benchmarks. Their compute-optimal curve
predicts that for any fixed problem difficulty, there exists an optimal
allocation between training compute (model size) and inference compute
(test-time search). Below that optimum, a smaller model with more inference
compute outperforms a larger model with less; above it, the larger model wins.

This framework is the closest theoretical grounding for the ``orchestration
economics'' claim, but it remains limited in three ways that motivate our study.
First, \citet{SnellEtAl2024} measure \emph{verifiable} reasoning tasks (math,
code) where a verifier can score intermediate steps; open-ended business
strategy has no verifier. Second, their test-time compute is
\emph{search-based} (sampling + scoring), whereas the multi-agent literature
uses \emph{structural} compute (builder--critic--refiner pipelines, debate
rounds) --- the two may have different cost--quality curves. Third, and most
fundamentally, their analysis is in units of FLOPs, not dollars: the
practitioner facing an API bill cannot pay in FLOPs. A dollar-matched empirical
test is missing.

\subsubsection{Self-refinement and the limits of iterative improvement}

The final reasoning structure relevant to this study is self-refinement: a
model critiquing and revising its own output. \citet{ShinnEtAl2023} introduced
Reflexion, in which a language agent maintains a reflective memory of past
failures and uses verbal reinforcement learning to improve subsequent trials.
On coding and decision-making tasks, Reflexion outperformed single-shot
baselines and even standard CoT by learning from its own errors.

Self-refinement raises a hypothesis directly tested in our Study~2: if a cheaper
model's initial output is below the frontier, can a low-cost refinement pass close
the gap? The multi-agent and self-refinement literatures imply that iterative
improvement should raise quality. But they do not tell us whether the
improvement is sufficient to reach parity with a frontier model that was not
itself refined --- the marginal comparison that determines substitution. Nor do
they tell us \emph{why} refinement succeeds or fails on open-ended tasks. Our
qualitative coding protocol (\cref{sec:qual-coding}) is designed to answer this
mechanistic question.

\subsection{LLM-as-a-judge evaluation methodology}

Because our outcomes are open-ended business solutions, evaluation itself is a
design object. \citet{ZhengEtAl2023} established LLM-as-a-judge as a scalable
evaluation paradigm, showing strong agreement between strong judge models and
human preferences on pairwise comparison tasks, while documenting three
systematic biases that threaten validity: position bias (preferring the
candidate presented first), verbosity bias (preferring longer outputs), and
self-enhancement bias (preferring outputs from the same model family as the
judge). Their mitigations --- swapping candidate order, controlled few-shot
examples, length normalization --- inform our blinding and order-shuffling
protocol.

The self-enhancement bias is particularly acute for multi-agent studies. If the
generator and judge share model architecture, the judge may reward the
generator's stylistic conventions rather than substantive quality.
\citet{VergaEtAl2024} extended the single-judge design into a jury: a Panel of
LLM evaluators (PoLL) composed of several smaller models from different
families correlated better with human judgments than a single large judge, at
lower cost, while mitigating the self-preference bias that afflicts same-family
evaluation. Our evaluation framework inherits this logic directly: three judges
from families disjoint from all generators, each under three frozen prompt
variants, scoring blind to cell and family. Because no judge shares a family with any generator, the ``persona penalty'' reported in \cref{sec:persona-penalty} cannot be attributed to same-family self-enhancement bias \citep{ZhengEtAl2023}.

\subsection{Open-science methods and economic measurement}
The fourth stream supplies the discipline this paper applies. \citet{NosekEtAl2018}
argue that preregistration distinguishes confirmatory from exploratory
analysis, converting transparency into credibility; our protocols, including
the four modification protocols (05--08), were frozen before the corresponding computations ran and are
timestamped in the OSF registration chain (\cref{tab:osf-registrations}), which
also discloses the two exceptions.
\citet{SteegenEtAl2016} show that conclusions can vary across reasonable
analytic choices and recommend reporting a multiverse of analyses; our
price-assumption panel (three blended input/output ratios) and pre-registered
robustness battery implement this idea. \citet{SandveEtAl2013} codify
reproducible computational research --- versioned scripts, recorded steps,
stable environments --- which our token ledger and frozen configuration files
follow.

Two statistical practices complete the foundation. \citet{AlbersLakens2018}
demonstrate that effect sizes estimated from pilots are unstable and bias
follow-up power analyses, so pilot data must never anchor confirmatory claims;
accordingly, our pilot data informed only feasibility gating, never the
confirmatory analysis. \citet{Lakens2013} makes the case for reporting effect sizes with
confidence intervals to support cumulative science; we report Welch's $d$ with
bootstrap intervals throughout. \citet{Cohen1960}'s $\kappa$ coefficient
provides our inter-rater and coder-agreement metric.

Finally, measuring the economics of inference requires separating quantity from
price. \citet{Diewert1976} formalized the decomposition of value changes into
price and quantity components in index-number theory; we apply the same logic
by storing tokens (quantity) and multiplying by a frozen price vector, so every
reported dollar is reproducible and re-priceable. \citet{TomlinsonEtAl2024}
quantify the environmental footprint of generative tasks, reinforcing that
inference cost --- economic and environmental --- is a first-order design
variable rather than an implementation detail.

\subsection{Integration: from persona economics to orchestration economics}

Our own prior work supplies the immediate empirical background. In a
cross-architecture evaluation of cultural personas across seven large language
models, persona effectiveness followed a contingency model rather than a
universal rule, and the Efficiency Gain Ratio (EGR) showed that small models
with personas achieved superior performance at 40--88\% lower cost
\citep{Hariri2026b}. The same line of work measured creativity with multi-judge
panels over thousands of generated business solutions, establishing Novelty as
the best-agreed creativity dimension \citep{Hariri2025a}, and showed that team
composition effects are contingent on task verifiability \citep{Hariri2025b}.

Integrating the four streams above with this lineage yields the gap and a
pre-registered hypothesis set spanning all three studies, each element grounded
in the literature or in our prior work. The Study~1 labels $H_1$--$H_5$ follow
the registered protocol exactly (OSF registration
\href{https://osf.io/mjb8x}{\texttt{osf.io/mjb8x}}, Protocol~01): $H_1$ is the
primary tier $\times$ structure interaction; $H_2$, $H_3$, and $H_4$ are the
tier, structure, and persona main effects; $H_5$ is the structure $\times$
persona interaction. Studies~2 and~3 add their own pre-registered hypotheses
below (Protocols~09 and~13).

\medskip
\noindent\textbf{Substitution hypothesis ($H_1$):} The test-time compute
framework \citep{SnellEtAl2024} motivates the prediction that structure can
substitute for model tier at the lower-tier end of the compute-optimal curve. We
test this prediction with a factorial tier $\times$ structure design, priced in
dollars, on open-ended tasks where no verifier exists. If the hypothesis holds,
team structure should raise quality more for cheaper models than for frontier
models.

\medskip
\noindent\textbf{Main-effect hypotheses ($H_2$, $H_3$):} The registered
protocol also specifies the tier main effect ($H_2$: frontier models score
higher than their lower-tier siblings, averaged over structure and persona)
and the structure main effect ($H_3$: teams score higher than solos, averaged
over tier and persona). We report both as secondary context for the economic
analysis.

\medskip
\noindent\textbf{Persona hypotheses ($H_4$, $H_5$):} Our prior work
\citep{Hariri2026b} showed that cultural personas shift quality within
single-model cells. $H_4$ extends the registered claim to this design:
persona-loaded conditions should score higher than no-persona conditions,
averaged over tier and structure. $H_5$ adds that the persona effect should be
larger for team structures than for solo structures (personas amplify
structural rigor). If personas primarily add stylistic variation without
substance, the persona-loaded cells should instead be Pareto-dominated by
their no-persona twins.

\medskip
\noindent\textbf{Refinement hypothesis (Study~2, Protocol~09):} The
self-refinement literature \citep{ShinnEtAl2023} suggests that a low-cost
critique-and-revise pass can recover quality gaps. Study~2 tests whether a
frozen, add-only refiner can close the residual gap between orchestrated
lower-tier outputs and frontier outputs --- and whether the mechanism is
additive coverage (more structure) or compressed insight (better density).
Study~2 was pre-registered as a separate study before any refinement run,
with its own feasibility gate (G1), superiority gate (G2), and three
directional criteria (H2.1--H2.3) frozen in the decision log (Protocol~09,
decisions D39--D40).

\medskip
\noindent\textbf{Authorship hypotheses (Study~3, Protocol~13):} Every pipeline
in the literature above --- and in Studies~1--2 --- is \emph{human-designed};
the registered-but-unexecuted A7 arm isolated authorship as the one design
variable left untested, and Study~3 executes it. Across the same five families,
the cheap model either authors and runs its own workflow (D-SELF) or runs one
authored by the frontier flagship (D-FRONT). The primary hypothesis is
\emph{directional}: frontier-authored pipelines outperform self-authored ones
(H-D1: D-FRONT $>$ D-SELF), on the ground that design capability transfers
through authored structure --- the same logic by which a cross-family refiner
beat within-family editors in the Study~2 extension. Two further hypotheses complete the set:
that model-authored pipelines are at least as good as the frozen human-designed
instrument (H-D2, two-sided), and that frontier-authored design stays cheaper
per quality point than frontier solo once design tokens are charged (H-D3,
economic). A pre-registered confirmatory replication (Protocol~16) then converts
the family-level effects into binding directional gates, testing whether the
cost--quality ordering transports to fresh problems.

\cref{sec:method} specifies the pre-registered design, apparatus, and analysis
plan that test this hypothesis set, with every run priced in dollars;
\cref{sec:results} reports the outcomes. One hypothesis is mechanistic rather
than purely economic: the qualitative coding protocol (\cref{sec:qual-coding})
tests the refinement hypothesis's prediction that orchestration buys
\emph{coverage} while the frontier wins on \emph{density}, by coding original,
refined, and frontier answers against a frozen codebook.



% ==================================================================
% 3  METHODOLOGY
% ==================================================================
\section{Methodology}\label{sec:method}

\subsection{Research Design}

The experiment is a $2 \times 2 \times 2$ between-cell factorial with model
\emph{family} as a blocking factor. The three manipulated factors are
\textbf{tier} (frontier flagship vs.\ lower-tier sibling within the same
generation), \textbf{structure} (solo single call vs.\ three-step team of
builder $\to$ critic $\to$ refiner), and \textbf{persona} (none vs.\ a fixed
cultural-persona mapping applied to each team role). Five model families ---
Kimi, Qwen, GPT-5.6, Claude, and Gemini --- supply both frontier and cheaper
tiers. Thirty business problems, stratified across knowledge-work domains,
are identical across all cells; no instruction wrappers are added to the
rendered prompt. Every design and statistical decision is time-stamped, numbered
(D1--D54), and registered in a frozen decision log; sixteen protocols ---
twelve standalone and four modifications of the original registration (05--08)
--- were registered before the corresponding computation ran, maintaining
the confirmatory/exploratory separation required by preregistration discipline
\citep{NosekEtAl2018}.

Two terms recur throughout the paper and are worth fixing precisely. A
\textbf{protocol} is a single pre-registered experiment plan: it states the
research question, the hypotheses, the models and problems to be run, the
analysis code, and the stopping rules --- and it is frozen (time-stamped and
registered on OSF) \emph{before} any data for that question are generated, so
that results cannot be retrofitted to expectations. An \textbf{arm} is one
experimental condition inside a protocol --- one specific combination of model
tier, pipeline structure, and persona setting that produces a batch of runs
(e.g.\ ``lower-tier team without persona'' is one arm). \cref{app:registry}
lists every protocol and arm with its role and status.

This paper reports three studies and two pre-registered replications.
\textbf{Study~1} is the factorial experiment above ($H_1$--$H_5$; Protocols
01--04 and their modification Protocols 05--08). \textbf{Study~2} asks whether a cheap
\emph{refinement pass} on an existing answer can close the remaining gap to
frontier quality (Protocol~09), followed by qualitative coding of what
orchestrated and frontier outputs actually contain (Protocol~10). A
pre-registered \textbf{Study~2 extension} (Protocol~11) repeats the refinement
test with two further families (GPT, Gemini) and adds a refiner-identity
contrast: each family refined by its own model versus a \emph{cross-family}
refiner editing another family's output. \textbf{Study~3}
manipulates pipeline design authorship with model-authored pipelines
(Protocol~13, \cref{sec:study3}). The two replications stress-test the core
claims out-of-sample: a confirmatory wave on 30 fresh problems across ten
families (Protocol~12, OSF~\href{https://osf.io/jg49x}{jg49x};
\cref{sec:wave}), and a transport test of Study~3's family-level effects on
new problems (Protocol~16; \cref{sec:study3conf}).

\begin{table}[ht]
  \centering
  \caption{OSF registration chain (Open-Ended Registrations branched from project
  node \texttt{8pwej}). The text cites Protocol numbers; this table is the
  authoritative GUID map.}
  \label{tab:osf-registrations}
  \footnotesize
  \begin{tabular}{@{}llll@{}}
    \toprule
    \textbf{Study / Protocol} & \textbf{OSF} & \textbf{Registered} & \textbf{Role} \\
    \midrule
    Study 1 (Protocols 01--04, A1/A2) & \href{https://osf.io/mjb8x}{mjb8x} & 2026-08-06 & factorial pre-registration$^{\dagger}$ \\
    Study 2 extension (Protocol 11)     & \href{https://osf.io/83c9u}{83c9u} & 2026-08-22 & refinement extension (GPT, Gemini) \\
    Wave (Protocol 12)                & \href{https://osf.io/jg49x}{jg49x} & 2026-08-23 & Study 1 confirmatory replication \\
    Study 3 (Protocol 13, v2)         & \href{https://osf.io/bwc93}{bwc93} & 2026-08-29 & authorship study (operative)$^{\ddagger}$ \\
    Study 3 (Protocol 13, G1)         & \href{https://osf.io/639eq}{639eq} & 2026-08-29 & feasibility-gate clarification \\
    Study 3 replication (Protocol 16) & \href{https://osf.io/zkm8v}{zkm8v} & 2026-08-30 & family-moderation transport test \\
    \bottomrule
  \end{tabular}
  \par\vspace{4pt}
  \begin{minipage}{0.98\textwidth}
    \footnotesize\raggedright
    $^{\dagger}$\,mjb8x \emph{memorializes} the Study~1 pre-registration: the
    protocol and decision log were frozen before data collection, but this OSF
    timestamp post-dates main-run generation (as the registration itself
    discloses).\\
    $^{\ddagger}$\,bwc93 is the operative Study~3 registration, superseding the v1
    registration (9ukjz) after a pre-analysis instrument audit; the 73 v1 runs
    were archived and never pooled.
  \end{minipage}
\end{table}

\textbf{Registered but unexecuted arms (disclosure).} The Study-1
pre-registration also specified four special mechanism arms --- A5 (a
critic-gated WHILE loop on the cheap team), A6 (a cost-matched frontier-solo
baseline), A7 (an adaptive per-problem pipeline designer), and A8 (a
shuffled persona-to-role mapping) --- each with a pre-registered
feasibility/futility gate. These gates were never evaluated and the arms were
never executed within Study~1: the main run was scoped to the eight cells of
the factorial design itself (2 tiers $\times$ 2 structures $\times$ 2 persona
settings, listed in \cref{app:registry}), and the Study-2 extension and the
replication
wave took budgetary priority. Per
no-file-drawer discipline we disclose their registered-but-unexecuted status
here rather than silently dropping them. The design-authorship question behind
A7 has since been executed as Study~3 (Protocol~13; \cref{sec:study3}); the
loop question behind A5 is the subject of a separate pre-registered protocol
(\cref{sec:future}).

\begin{figure}[ht]
  \centering
  \includegraphics[width=\textwidth]{fig1_design_schematic.png}
  \caption{Experimental Design: $2 \times 2 \times 2$ Factorial + Apparatus.}
  \label{fig:design}
\end{figure}

\subsection{Models Evaluated}\label{sec:models}

\subsubsection{Frontier Models (Five Families)}

Each family contributes its current \emph{mainline} flagship, pinned on
2026-07-30: \texttt{kimi-k3}, \texttt{qwen3.7-max-2026-06-08},
\texttt{gpt-5.6-sol}, \texttt{claude-opus-5}, and
\texttt{gemini-3.1-pro-preview}. Above-mainline tiers (e.g.\
\texttt{claude-fable-5}) are deliberately excluded because they would break
cross-family comparability. Gemini's preview status carries a pre-registered
fallback to \texttt{gemini-2.5-pro} if the model is retired mid-experiment; the
fallback was never invoked. All models run with generator temperature $= 1$, a
uniform value forced by the OpenAI and Moonshot APIs (which reject other values)
and applied identically across families for \emph{ceteris paribus}
comparability~(D27).

\subsubsection{Lower-Tier Siblings}

\begin{sloppypar}
Each family contributes one lower-tier model from the same generation:
\texttt{kimi-k2.6}, \texttt{qwen3.7-plus-2026-05-26}, \texttt{gpt-5.6-luna},
\texttt{claude-haiku-4-5-20251001}, and \texttt{gemini-3.1-flash-lite}. Price
ratios (frontier : cheaper) range from 3.1$\times$ to 8$\times$ across families.
The size of this price gap differs from family to family; our statistical model
simply gives each family its own baseline (a \emph{family random effect},
\cref{sec:inferential}), so these differences in price gap cannot distort the
main tier $\times$ structure comparison.
\end{sloppypar}

\subsubsection{Replication-Wave Models (Five Additional Families)}

The confirmatory wave (Protocol~12, \cref{sec:wave}) repeats the three core
cells (C-S-N, C-T-N, F-S-N; persona excluded by design, D50) on 30 fresh
problems and adds five families beyond the core set, pinned on
2026-08-22/23 per Protocol~12 \S3: Thinking Machines
(\texttt{Inkling} / \texttt{Inkling-Small}), OpenAI open-weight
(\texttt{gpt-oss-120b} / \texttt{gpt-oss-20b}), Google open-weight
(\texttt{gemma-4-31B-it} / \texttt{gemma-3n-E4B-it}), Mistral
(\texttt{mistral-medium-2604} / \texttt{mistral-small-2603}), and MiniMax
(\texttt{MiniMax-M3} / \texttt{MiniMax-M2.7}). The five core families run their
pinned Study-1 models unchanged, with one disclosed routing-only deviation:
the Kimi lane runs on the native Moonshot provider because the Together-hosted
\texttt{kimi-k2.6} was retired before the wave (D48; identical prompts and
decoding, provider route auditable in the ledger).

\subsection{Orchestration Structures and Personas}

\textbf{Solo} runs invoke a single model call on the rendered problem.
\textbf{Team} runs execute a three-step pipeline --- builder produces a draft,
critic critiques it, refiner produces the final answer --- with all intermediate
outputs discarded from the judged answer (and therefore from the quality score),
while the cost ledger records every step: the per-run token total is the sum of
the builder, critic, and refiner calls, so team costs include the full pipeline,
not just the final step. The one gap is retry accounting: when a step failed
and was retried, the failed attempt itself logged zero tokens and zero latency
(an instant transport rejection, consuming nothing billable), but on 110 of the
950 completed team runs the ledger's run total then omitted tokens from
\emph{completed} earlier steps, so recorded cost understates true cost on those
runs; we bound this in \cref{sec:limitations}. Persona is implemented as a fixed mapping from role
(builder/critic/refiner) to a cultural persona description, identical across all
problems and families \citep{Hariri2026b}. The mapping is pre-registered and
never varied per-problem; it is therefore a treatment level, not a confound.

\subsection{Stimuli and Response Generation}

Thirty domain-stratified business problems are used. Fifteen are inherited
verbatim from a prior validated set \citep{Hariri2025a}; the remaining fifteen
were drafted for this study so that the full grid covers six domains with five
problems apiece (the post-cutoff contamination-control cases belong to the
replication wave's fresh problem set, \cref{sec:wave}). Each problem is
rendered as \texttt{background\_context + "\textbackslash n\textbackslash n" +
business\_problem} with no instructional wrapper --- adding wrapper text would
introduce an unmeasured prompt treatment. Generator temperature is fixed at~1
(\cref{sec:models}). Every run is time-stamped and flagged for whether it
executed inside the pre-registered preferred window (14:00--00:00 UTC), which
captures the Qwen off-peak discount; window-free providers (Kimi, Together-hosted
models, Western providers) are time-invariant and therefore unaffected.
DeepSeek pricing was stable throughout the experiment (no peak-surge window was
in effect during the execution period).

\subsection{Evaluation Framework}\label{sec:eval}

\subsubsection{Judge Model Selection}

Three judges from families disjoint from every generator: \texttt{deepseek-chat}
(resolved to \texttt{deepseek-v4-flash}), \texttt{zai-org/GLM-5.2}, and
\texttt{meta-llama/Llama-3.3-70B-Instruct-Turbo}. The disjointness rule~(D33)
guards against self-enhancement bias \citep{ZhengEtAl2023}: even when deepseek
is the registered fallback generator, the fallback was never invoked, so
disjointness holds in the executed data.

\subsubsection{Scoring Rubric}

Five 1--5 dimensions are scored: \textbf{Novelty, Usefulness, Flexibility,
Elaboration, and Cultural Sensitivity}. The composite
primary outcome is the mean of the five dimensions. Three prompt variants
are used per judge, frozen verbatim and never re-optimized after a prior
judge-prompt optimization attempt collapsed $\kappa$ in earlier work
\citep{Hariri2025a}. Every solution is scored by all three judges under all
three variants (nine scores per solution), with candidate order shuffled per
variant to neutralize position bias \citep{ZhengEtAl2023}. Judges run at
temperature~0 for determinism.

\subsubsection{Inter-Rater Reliability}

Pairwise judge correlations on per-run composite scores are 0.87
(deepseek--zhipu), 0.61 (deepseek--meta), and 0.49 (meta--zhipu). The
panel-level reliability exceeds any single judge's correlation with the latent
quality; this mirrors the jury logic of \citet{VergaEtAl2024} in which diversity
of smaller judges outperforms a single large judge.

\subsection{Statistical Analysis}

\subsubsection{Effect Sizes}

Welch's $d$ with bootstrap 95\% confidence intervals (5{,}000 resamples, seed
42) throughout \citep{Lakens2013}. Estimation-first reporting: intervals, not
$p$-values, carry the weight of the argument; $p$-values are reported for
completeness and to comply with sequential-analysis boundaries
(\cref{sec:inferential}).

\subsubsection{Inferential Statistics}\label{sec:inferential}

The primary test of $H_1$ is the pre-registered interaction contrast
$d = (FT - FS) - (CT - CS)$, expected negative (team structure helps cheaper
models more than frontier models). Inference uses a mixed model with random
effects for \emph{problem} and \emph{family}, solving the pseudo-replication
identified in the author's
prior work \citep{Hariri2026b}. Sequential analysis with an O'Brien--Fleming
spending function (protocol~03) governs interim looks and the final stopping
boundary; the Bayesian JZS Bayes factor is reported as recommended and
pre-registered in the robustness battery\footnote{Protocol 06 --- Extended
Interim-2 Analysis Methods, items 6.1--6.7; OSF registration
\url{https://osf.io/mjb8x}.} but
does not gate the decision. Family-wise corrections within secondary hypotheses
use Benjamini--Hochberg FDR ($\alpha = 0.05$); Holm is used for the small set of
directional family comparisons in Study~2.

Later studies and replications follow the same estimation-first discipline with
study-specific designs. The \textbf{Study~2 extension} (Protocol~11) uses paired
within-problem tests with Holm adjustment across arms, plus a two-sided
refiner-identity contrast (no pre-committed direction). The \textbf{wave}
(Protocol~12) fits an LMM with random intercepts for family and variance
components for problem and family $\times$ problem, evaluated against binding
gates W1--W7; when the dual-variance-component model was numerically singular, a
pre-declared fallback substituted a parsimonious structure with bit-identical
fixed effects (D52). \textbf{Study~3} (Protocol~13) tests the authorship
contrast with paired $t$-tests within family $\times$ problem, Holm-adjusted
across the five families, Novelty as the primary dimension and the composite as
a robustness variant (conclusions are DV-invariant). \textbf{Protocol~16}
converts Study~3's point estimates into binding directional gates: problem-level
Welch tests of D-SELF against the frozen wave F-S-N baseline, one-sided $p$ by
normal approximation, and a TOST equivalence gate for GPT-5.6 (90\%~CI within
$\pm 0.30$), executed by a hash-locked analysis script frozen before judging.

\subsubsection{Economic Analysis}

Cost is stored as tokens in a run-level ledger and re-priced through a frozen
price vector (\texttt{config/price\_vector.json}, version 2026-07-30). Because
the ledger does not split input from output tokens, cost is computed under a
blended price $\mathrm{blend}(\phiBlend) = \phiBlend \cdot p_{\mathrm{out}} +
(1 - \phiBlend) \cdot p_{\mathrm{in}}$ with a primary $\phiBlend = 0.50$ and a
sensitivity panel $\phiBlend \in \{0.25, 0.50, 0.75\}$ (\citealt{Diewert1976}
provides the price--quantity decomposition precedent; \citealt{SteegenEtAl2016}
the multiverse-robustness precedent). Frontier-membership claims must be stable
across all three $\phiBlend$. Studies~3 and the P16 confirmatory use a newer
ledger that records prompt and completion tokens separately for every call, so
their costs are priced exactly at provider list rates; the
blended-$\phiBlend$ treatment applies to Studies~1--2 and the wave, where only
total tokens were logged. The \textbf{Tiered Efficiency Gain Ratio (T-EGR)}
is $(\mathrm{Quality}_{\mathrm{cell}} / \mathrm{Quality}_{F\text{-}S\text{-}N})
/ (\mathrm{Cost}_{\mathrm{cell}} / \mathrm{Cost}_{F\text{-}S\text{-}N})$,
computed within each family and pooled; F-S-N is the baseline. \textbf{ICER} is
reported relative to the lowest-cost cell. Judge cost is excluded from unit T-EGR
(it is shared infrastructure, not cell compute; a disclosed exclusion).

\subsubsection{Robustness Checks}

The pre-registered robustness battery for Study~1 (Protocol~06,
items 6.1--6.7) comprises
bootstrap confidence intervals on all primary contrasts (6.1),
Benjamini--Hochberg FDR on secondary contrast families (6.2), a Type-II
factorial ANOVA sensitivity layer with partial $\eta^2$ (6.3), JZS Bayes
factors for non-significant contrasts (6.4), leave-one-family-out
sensitivity (6.5), leave-one-problem-out generalizability (6.6), and a
sensitivity power analysis in place of post-hoc power (6.7); all seven are
reported in \cref{sec:robustness-families}. After the data
lock, an explicitly exploratory modification (Protocol~07, registered in the
decision log but not part of the OSF package) added five analyses
(7.1--7.5); the two
interpreted in this paper are the random-effects meta-analysis across the five
Study-1 families (7.4, \cref{sec:robustness-families}) and the persona
dimension decomposition (7.5, \cref{sec:persona-penalty}); the remaining
three are archived in the replication package. All Protocol-07 results are labeled
exploratory wherever reported. For Study~1, the item-7.4 meta-analysis confirms
the tier $\times$ structure interaction generalizes across families
(pooled $d = -0.270$, $I^2 = 28.5\%$). Studies~2 and~3, the replication
wave, and the P16 confirmatory carry their own pre-registered analysis gates,
described with their respective designs
(\cref{sec:wave,sec:study2,sec:study3,sec:study3conf}).

\subsubsection{Qualitative Coding Protocol (Study~2)}\label{sec:qual-coding}

Thirty triples (original CTP / refined CTP / same-problem same-family FSN) are
sampled deterministically (seed 20260801, 10 per family). Two blinded LLM coders
perform a two-wave coding: free induction on a subset, then closed coding
against a frozen 13-code codebook.\footnote{Codebook development followed the
pre-registered qualitative analysis protocol (OSF registration
\url{https://osf.io/mjb8x}, file \texttt{02-QUALITATIVE-ANALYSIS-PROTOCOL.md}).} The codebook is organized
around three research questions: \textbf{RQ1} asks whether orchestration adds
\emph{insight} (compressed reasoning, named concepts, diagnostic reframing) ---
six codes; \textbf{RQ2} asks whether orchestration adds \emph{structure}
(scaffolding, timelines, matrices) --- six codes; \textbf{RQ3} asks whether
orchestration adds \emph{bloat} (redundancy) --- one code. Cohen's $\kappa$ gate is 0.70; codes below threshold
are downgraded to medium reliability and disclosed \citep{Cohen1960}. The prompt
presented to coders contains the three texts only --- no cell label, family, or
score, preserving blind evaluation.

\subsection{Validity and Reliability}

\textbf{Construct validity:} the rubric separates Novelty from
Usefulness/Elaboration, so a stylistic shift (e.g.\ more discourse markers from
a persona) is not conflated with substantive quality. \textbf{Internal validity:}
frozen configurations, blinded judging, shuffled order, disjoint judge families.
\textbf{External validity:} single task domain (English business strategy) is
the acknowledged boundary. \textbf{Disclosed threats.} (i)~\emph{Retry cost accounting.} On 110 of the 950
completed team runs, a retried step caused the ledger's run total to omit
tokens from completed earlier steps, so reported team costs understate true
cost on those runs (median true-to-recorded ratio 1.69; about 6\% of team-cell
cost in aggregate). The bias favors the team cells and therefore favors the
paper's central claim. It is disclosed but not corrected, since correcting it
could only shrink, not reverse, the cost advantage. (ii)~\emph{Truncation at source.} A small set of
very long runs was capped by the transport layer before storage; each run
records whether this happened, and the affected runs are audited separately
rather than silently excluded.

\subsection{Code and Data Availability}
All analysis code, the frozen token ledger, judge rationales, the decision log
(D1--D54), and Protocols~09--10 and~13--16 are committed to the
\texttt{orchestration-economics} repository. The pre-registration trail is the
OSF registration chain mapped in \cref{tab:osf-registrations}: each OSF
registration archives its protocol file, and the table discloses the two
exceptions --- the Study~1 memorialization (mjb8x) and Study~2
(Protocols~09--10), which are committed to the repository and frozen in the
decision log (D39) but were not separately OSF-registered.
A materialized results layer (D54) ships every statistical test as a flat CSV
--- per-solution judge scores for all studies, the wave gate statistics, the
heterogeneity batteries, and the qualitative codebook with per-code $\kappa$
--- so that no number in this paper requires re-running code to inspect.
A reproduction runbook
\citep{SandveEtAl2013} lists the commands that regenerate every number reported
in~\cref{sec:results}.

\subsection{AI-Assisted Research Disclosure}\label{sec:ai-disclosure}

This study used large language models as research tools under continuous human
supervision. Specifically, \textbf{Kimi} (Moonshot AI, \texttt{kimi-2.6 and kimi-k3}) and
\textbf{Qwen} (Alibaba Cloud, \texttt{qwen3.7-max and qwen3.8-max}) were employed in the
following capacities: (1)~\textbf{ideation} --- brainstorming experimental
design alternatives, hypothesis formulations, and analysis strategies during the
protocol-drafting phase; (2)~\textbf{code generation} --- producing initial
versions of the execution pipeline, statistical analysis scripts, and
figure-generation code, all of which were reviewed, corrected, and approved by
the author before execution; (3)~\textbf{audit} --- reviewing citation accuracy,
checking statistical assumptions, and identifying inconsistencies across
protocol documents; (4)~\textbf{execution} --- running the generation and
judging campaigns through API calls, with all configurations frozen and logged
before each run. Higher-tier models were employed for ideation, audits, and
task execution; lower-tier models were employed for code generation.

No LLM made unsupervised decisions about the experimental design, statistical
analysis, or the interpretation of results. All twelve pre-registered protocols,
all 54 design decisions (D1--D54), and all analysis choices were authored,
reviewed, and approved by the human author. Outputs were
validated against expected ranges, and every reported number traces to a frozen,
version-controlled analysis pipeline. The qualitative coding protocol
(\cref{sec:qual-coding}) used two LLM coders with blinded evaluation --- a
design choice that converts the LLM's systematic biases into a measured
inter-rater reliability coefficient ($\kappa = 0.824$) rather than an unmeasured
confound.

A concrete oversight episode illustrates why continuous human audit is not
procedural decoration. When the AI agent implemented the Study~3
design-authorship instrument, it silently introduced two deviations from the
product feature under study (DeeWaanAI's auto-designer): it hard-capped
pipelines at four steps (the platform allows up to seven) and wrote a bespoke
designer prompt, whereas the platform uses a fixed internal prompt with
free-form, task-adapted roles. Both deviations were invisible in the code's
behaviour --- every pilot design complied, so no error surfaced --- and were
caught only when the human author audited the instrument's rationale against
the platform source code before confirmatory data collection. The instrument
was corrected and the correction registered \emph{before} the Study~3
campaign (D56). The lesson for researchers delegating to AI agents: an AI
assistant's silent defaults are themselves experimental-design decisions, and
they diverge from the intended construct precisely where the researcher did
not specify --- audit the rationale behind every parameter, not only the
parameters named in the protocol.

The use of LLMs as evaluation judges (\cref{sec:eval}) is a distinct
methodological choice from the use of LLMs as research tools: the judges
evaluate \emph{generated business solutions}, not the research process itself.
The disjointness rule~(D33) ensures that no judge shares architecture with any
generator, mitigating the self-enhancement bias documented by
\citet{ZhengEtAl2023}.

\textbf{Platform use and conflict-of-interest disclosure.} All pipeline
generation and execution in Studies 1--3 ran on \textbf{DeeWaanAI}
(\url{https://deewaanai.com}), a commercial orchestration platform; the
auto-designer evaluated in Study~3 and Protocol~16 is DeeWaanAI's production
feature, used unmodified (fixed internal design prompt, free-form roles,
up to seven steps with critic-gated loops). The author is affiliated with
DeeWaanAI, and the findings directly inform the platform's routing and
persona features (\cref{sec:conclusion}). This relationship is disclosed so
readers can weight the translational claims accordingly; all protocols,
gate outcomes (including failed ones), and attrition events are
pre-registered and reported as-is, and the full data and code are publicly
available.

% ==================================================================
% 4  RESULTS
% ==================================================================
\section{Results}\label{sec:results}

This section reports three studies that escalate a single business question ---
\emph{orchestrate cheaper models, or pay for the frontier?} \textbf{Study~1}
(\cref{sec:study1}) tests whether pipeline structure substitutes for model tier
and prices the answer in dollars. \textbf{Study~2} (\cref{sec:study2}) explains
\emph{why} that substitution stops short of the frontier. \textbf{Study~3}
(\cref{sec:study3}) asks whether it matters \emph{who designs} the pipeline.
Each study's headline claims are then stress-tested out-of-sample in a
pre-registered replication, so the rules we recommend are the ones that survived
contact with new data --- and the ones that did not are reported as failures,
not buried.

\subsection{Study 1: Does Structure Substitute for Model Tier?}\label{sec:study1}

\subsubsection{The interaction effect: structure is a lower-tier
amplifier}\label{sec:interaction}

From a resource-allocation perspective, the most important finding is that
pipeline structure delivers higher marginal returns at the lower tier than at
the frontier tier. The pre-registered interaction contrast is $d = -0.227$
(bootstrap 95\% CI $[-0.410, -0.040]$, Welch $t = -2.45$, $p = 0.0145$), with
the sequential-analysis boundary crossed at the scheduled final look. For a
firm making an AI budget decision, this implies that pipeline dollars are better
spent upgrading lower-tier models to team configurations than adding team overhead to
already-expensive frontier models. The JZS Bayes factor for the interaction is
1.002 (inconclusive), reported as recommended but not gating the decision.

The secondary effects provide additional strategic context. The \textbf{tier
main effect} ($H_2$: $d = +0.294$, CI $[+0.201, +0.382]$, $\BF_{10} = 1.7 \times
10^7$) confirms that frontier models command a quality premium, but the
magnitude is dwarfed by the cost difference. The \textbf{structure main effect}
($H_3$: $d = +0.117$, CI $[+0.030, +0.205]$) shows that teams slightly outperform
solos on average --- a small but consistent productivity gain from
orchestration. The \textbf{persona main effect} ($H_4$) is large and in the
\emph{opposite} direction to the pre-registered expectation: $d = -0.691$ (CI
$[-0.761, -0.620]$, $\BF_{10} = 1.0 \times 10^{50}$). Persona-loaded cells
score \emph{lower} than their no-persona twins, representing a failed
investment in stylistic differentiation. From a business perspective, personas
degrade output quality sharply while reducing cost only marginally (persona
cells emit fewer tokens) --- a value-destroying configuration. An exploratory dimension decomposition (Protocol~07, item
7.5) shows that Novelty is the sole
dimension where personas help (+0.154, $p_{\mathrm{FDR}} = 0.001$) while
Usefulness ($-1.024$) and Elaboration ($-0.966$) drive the reversal. The
structure $\times$ persona interaction ($H_5$) is not significant ($F = 1.23$, $p =
0.27$).

\begin{quote}
\textbf{Strategic Implication:} Pipeline structure is a lower-tier amplifier,
not a universal substitute: the marginal return to team structure is higher at
the lower tier, although the highest absolute return remains the lower-tier
solo call (\cref{sec:pareto}). Personas, by
contrast, represent value destruction --- they degrade quality while reducing
cost only marginally.
\end{quote}

\begin{figure}[ht]
  \centering
  \includegraphics[width=\textwidth]{fig21_sequential_boundary.png}
  \caption{Sequential analysis with O'Brien--Fleming spending function. The test
    statistic $z = -2.447$ crosses the interim-2 boundary $|z| \ge 1.977$
    (dashed line). The shaded region is the continuation region. The
    pre-registered decision is to stop and report $H_1$ as supported.}
  \label{fig:sequential}
\end{figure}

\begin{figure}[ht]
  \centering
  \includegraphics[width=\textwidth]{fig2_forest_plot.png}
  \caption{Quality by Experimental Cell with 95\% Confidence Intervals.}
  \label{fig:forest}
\end{figure}

\subsubsection{The persona penalty: style without substance}\label{sec:persona-penalty}

Every persona-loaded cell except the budget-extreme cheap solo is
Pareto-dominated on the cost--quality plane (\cref{sec:pareto}): C-T-P, F-S-P,
and F-T-P are each beaten on quality by lower-cost no-persona cells. Personas
do \emph{not} raise cost --- persona cells emit fewer tokens and price
slightly \emph{below} their no-persona twins --- but the small saving buys a
large quality loss, so the twin comparison on quality is uniformly negative
($d = -0.691$). The lone survivor, C-S-P, remains on the frontier only because
it is the cheapest configuration tested; its quality is the lowest of all
eight cells. Its headline efficiency ratio ($\TEGR$~8.60, the highest in the
study) is purely a small-denominator artifact of that minimal cost --- it
marks the budget floor, not a recommended configuration, since the same tier
without a persona (C-S-N) buys 0.30 points of quality for an additional
\$0.005. The penalty is largest on dimensions that rubric readers
associate with substance (Usefulness, Elaboration), smallest on the one most
susceptible to stylistic signaling (Novelty). This pattern mirrors the
contingency finding in the author's prior work \citep{Hariri2026b} and is
consistent with the jagged-frontier mechanism \citep{DellAcquaEtAl2023}:
personas push the generator toward a more stylized region of its output space
that the judge, correctly, penalizes on substance.

\begin{quote}
\textbf{Strategic Implication:} Personas degrade output quality sharply while
reducing cost only marginally --- they are value-destroying for
knowledge-work tasks. Firms should treat
them as a failed investment rather than a quality lever.
\end{quote}

\begin{figure}[ht]
  \centering
  \includegraphics[width=0.85\textwidth]{fig22_persona_radar.png}
  \caption{Quality-dimension radar for persona-loaded (red, dashed) versus
    no-persona (blue, solid) cells. The persona polygon is visibly smaller on
    Usefulness and Elaboration, slightly larger on Novelty --- the visual
    signature of the style-over-substance tradeoff.}
  \label{fig:persona-radar}
\end{figure}

\begin{figure}[ht]
  \centering
  \includegraphics[width=\textwidth]{fig10_persona_dimensions.png}
  \caption{Persona Effect by Quality Dimension (FDR-Corrected).}
  \label{fig:persona-dims}
\end{figure}

\subsubsection{The frontier gap: lower-tier teams do not beat frontier solo}\label{sec:frontier-gap}

Paired comparisons between cheaper-team-with-persona (CTP) and the
frontier-solo-no-persona baseline (F-S-N) are negative in all five families
(kimi $d = -0.71$, claude $d = -1.42$, qwen $d = -2.17$; all CIs exclude zero
after Holm correction). This residual gap is the empirical target that
Study~2 (\cref{sec:study2}) asks a cheap refinement pass to close.

\begin{quote}
\textbf{Strategic Implication:} As executed, lower-tier teams do not beat a single
frontier call on matched problems. The substitution thesis is bounded: structure
improves lower-tier models but cannot fully close the frontier gap without additional
investment.
\end{quote}

\subsubsection{The Pareto frontier: optimal configurations for budget-constrained firms}\label{sec:pareto}

The cost--quality map is shown in \cref{fig:cost-scatter}. At the primary
$\phiBlend = 0.50$, the cell means are:

\begin{table}[ht]
  \centering
  \caption{Cell summary statistics (primary analysis, $\phiBlend = 0.50$).}
  \label{tab:cell-summary}
  \begin{tabular}{@{}lrrrr@{}}
    \toprule
    \textbf{Cell} & \textbf{$n$} & \textbf{Cost (USD)} & \textbf{Quality} & \textbf{T-EGR} \\
    \midrule
    C-S-P & 285 & 0.0145 & 4.210 & 8.60 \\
    C-S-N & 285 & 0.0194 & 4.507 & 6.15 \\
    C-T-P & 285 & 0.0715 & 4.327 & 1.58 \\
    C-T-N & 285 & 0.0734 & 4.531 & 1.49 \\
    F-S-P & 190 & 0.0670 & 4.406 & 1.36 \\
    F-S-N & 190 & 0.0835 & 4.586 & 1.00 (base) \\
    F-T-N & 190 & 0.2951 & 4.610 & 0.27 \\
    F-T-P & 190 & 0.3044 & 4.370 & 0.27 \\
    \bottomrule
  \end{tabular}
\end{table}

\begin{figure}[ht]
  \centering
  \includegraphics[width=\textwidth]{fig3_cost_scatter.png}
  \caption{Cost-Effectiveness Frontier: 8 Cells + Study-2 Refined CTP.}
  \label{fig:cost-scatter}
\end{figure}

\begin{figure}[ht]
  \centering
  \includegraphics[width=\textwidth]{fig6_tegr_bars.png}
  \caption{Tiered Efficiency Gain Ratio (T-EGR) by Cell.}
  \label{fig:tegr-bars}
\end{figure}

The Pareto frontier is \textbf{identical across all three $\phiBlend$}:
\{C-S-P, C-S-N, C-T-N, F-S-N, F-T-N\}. The dominated set \{C-T-P, F-S-P,
F-T-P\} contains every persona-loaded team cell --- the economic echo of the
\cref{sec:persona-penalty} persona penalty. Frontier membership is stable across
the sensitivity panel, satisfying the pre-registered claim rule.

\begin{figure}[ht]
  \centering
  \includegraphics[width=\textwidth]{fig19_phi_sensitivity.png}
  \caption{Pareto frontier stability across the price-assumption sensitivity
    panel. Left: $\phiBlend = 0.25$ (input-heavy blend). Center: $\phiBlend =
    0.50$ (primary analysis). Right: $\phiBlend = 0.75$ (output-heavy blend).
    The frontier set \{C-S-P, C-S-N, C-T-N, F-S-N, F-T-N\} is identical in all
    three panels, satisfying the pre-registered cross-$\phiBlend$ stability
    rule.}
  \label{fig:phi-sensitivity}
\end{figure}

ICER (vs.\ the lowest-cost cell, C-S-P): C-S-N is the best buy at +\$0.016 per
point of composite quality; F-T-P is the worst at +\$1.81 per point. The
refined kimi C-T-P point (not in the table) sits at \$0.0654 / quality~4.585 ---
absolute parity with pooled F-S-N (4.586) at 78\% of its cost, with the caveat
that the paired within-family comparison remains negative (disclosure
per~\cref{sec:refinement}).

\begin{quote}
\textbf{Strategic Implication:} The economically optimal strategy is tier
substitution at the lower-tier end: a single call from a cheaper model
(C-S-N) delivers 98\% of frontier quality at 23\% of the cost, and the
three-step team variant (C-T-N) buys coverage at 88\% of the cost. The worst
dollar in the experiment is spent on frontier teams --- they cost
3.5$\times$ the frontier-solo price for negligible quality gain.
\end{quote}

\subsubsection{Robustness across model families and problems}\label{sec:robustness-families}

The pre-registered robustness battery (Protocol~06) supports the
\cref{sec:interaction} decision.
Leave-one-family-out (item 6.5): $d$ ranges from $-0.143$ (drop kimi) to $-0.303$ (drop
gemini), sign stable in all five folds; no single family drives the
interaction. Leave-one-problem-out (item 6.6): $d$ ranges from $-0.212$ to $-0.247$ across
all 19 folds; no problem drives the interaction. Sensitivity power at the
smallest cell ($n = 190$) detects $d = 0.288$ at 80\% power; the observed $|d|
= 0.227$ is detected through the within-subjects contrast over all 1{,}900 runs
rather than per-cell power~(item 6.7). The Type-II ANOVA sensitivity layer (item 6.3)
agrees: tier $\times$ structure interaction $F = 6.90$, $p = 0.0087$, partial
$\eta^2 = 0.0036$; residuals are non-normal, so the bootstrap intervals and
Welch degrees of freedom remain the authoritative inference, as pre-specified.
Length adjustment on the $H_1$ contrast
yields $p = 0.63$ --- the interaction is not a length confound. The
random-effects meta-analysis over the five families (exploratory, Protocol~07,
item 7.4) gives a pooled $d =
-0.270$ (CI $[-0.472, -0.069]$, $z = -2.64$, $p = 0.008$) with low-moderate
heterogeneity ($\tau^2 = 0.015$, $I^2 = 28.5\%$), confirming the interaction
generalizes.

Across all battery perturbations the interaction keeps its sign and the
decision is unchanged. The leave-one-out sensitivity (\cref{fig:loo})
and random-effects meta-analysis (\cref{fig:meta-analysis}) confirm that no
single family, problem, or outlier drives the result.

\begin{quote}
\textbf{Strategic Implication:} The findings are robust to judge, model family,
problem domain, and output-length variations. The interaction is small but
reliable and generalizes across the five leading LLM families --- a stable
result for firms making platform-agnostic investment decisions.
\end{quote}

\begin{figure}[ht]
  \centering
  \includegraphics[width=\textwidth]{fig8_robustness_loo.png}
  \caption{Robustness: Leave-One-Out Sensitivity (A2 Battery).}
  \label{fig:loo}
\end{figure}

\begin{figure}[ht]
  \centering
  \includegraphics[width=\textwidth]{fig9_meta_analysis_forest.png}
  \caption{Random-Effects Meta-Analysis: Interaction by Family.}
  \label{fig:meta-analysis}
\end{figure}

\subsubsection{Exploratory heterogeneity: domain contingency, a capability-gap mechanism, and the persona decomposition}\label{sec:exploratory-heterogeneity}

\textbf{Status declaration.} Everything in this subsection is \emph{exploratory
and post-hoc}: it was formulated after data lock, is not part of the
pre-registered hypothesis set $H_1$--$H_6$, and consists exclusively of
re-analyses of the frozen Study-1 artifacts (no new generation or judging
calls). We report it because it generates concrete, falsifiable predictions ---
which were then tested by the pre-registered confirmatory replication wave
(\cref{sec:wave}). Anticipating that section: \textbf{all three upgrade
candidates below (domain contingency, the capability-gap mechanism, the
marketing variance advantage) FAILED their binding out-of-sample gates} and
therefore remain exploratory; the numbers in this subsection stand as the
Study-1 record, and the wave numbers are reported alongside in
\cref{sec:wave}~(Nosek et~al., 2018; Open Science Collaboration, 2015).

\textbf{Where does heterogeneity live?} Three screening analyses localize the
variation. Problem difficulty does not moderate the team effect
($r = 0.153$, $p = 0.53$ across the 19 problems). Family quality rankings are
remarkably stable across problem domains (Kendall's $W = 0.922$;
Kendall \& Babington Smith, 1939), so family choice is largely
domain-invariant. But measurement reliability itself is domain-heterogeneous:
single-judge agreement ICC(3,1)~(Shrout \& Fleiss, 1979) ranges from $0.235$
(HR/organization) to $0.523$ (marketing), with panel reliability ICC(3,9)
between $0.73$ and $0.91$ --- HR answers are simply harder for any judge,
human or LLM, to score consistently.

\textbf{The team effect is domain-contingent.} Pooling the cheap tier, the
team-minus-solo quality effect varies from $+0.004$ (marketing --- the only
domain where teams match solo) through strategy ($-0.036$), HR ($-0.053$),
operations ($-0.058$), product innovation ($-0.103$), to finance
($-0.106$). In marketing the team advantage extends beyond the mean: team
output variance is significantly \emph{lower} than solo (Brown--Forsythe test,
$p = 0.020$; Brown \& Forsythe, 1974), and downside risk shrinks across cells
($P(\text{composite} < 4.0)$: solo $2.8\%$, team $1.8\%$, frontier $1.1\%$).
The ensemble-variance logic is classical --- combining models reduces error
when members are accurate and diverse~(Hansen \& Salamon, 1990; Dietterich,
2000) --- but here it appears only for communication-intensive tasks.

\textbf{A capability-gap mechanism.} The family-level team effect (team minus
frontier solo) correlates almost perfectly with the family's capability gap
(frontier minus cheap solo)
(Pearson $r = -0.978$, $p = 0.004$; Spearman $\rho = -1.000$; $n = 5$
families): the team effect is most positive where the cheap model already
matches its frontier sibling (kimi, $-0.164$ gap --- the cheap model scores
\emph{above} its frontier --- and $+0.105$ team effect) and most negative
where the cheap model lags farthest below it (claude, $+0.364$ gap,
$-0.292$ team effect). One caution applies beyond the small $n$: the
correlation is partly mechanical, because the team effect and the gap both
contain the frontier mean. With only five families
this remains a hypothesis generator, not a test; the replication wave expands
the family universe specifically to power this correlation properly
(moderator analysis logic as in Borenstein et~al., 2009).

\textbf{Length control reveals an asymmetry.} A length-adjusted mixed model
(composite $\sim$ log-tokens $+$ cell, random family and problem intercepts;
Westfall et~al., 2014) leaves the frontier advantage intact ($+0.078$,
$p < 0.001$) but reduces the team's small raw edge ($+0.026$) to $-0.018$
($p = 0.40$). We deliberately do \emph{not} read this as mediation: quality is
judged on visible full text, and team structure mechanically produces longer
visible answers, so controlling length partially controls the intervention
itself. The defensible reading is the asymmetry --- frontier superiority
survives length adjustment; the team edge does not. The same model locates the
remaining economics: frontier solo is \emph{faster} than cheap teams ($94.8$s
vs.\ $131.7$s) and cheap solo dominates efficiency ($295.8$ vs.\ $28.1$
quality points per dollar-minute for teams).

\textbf{Persona decomposition.} Extending the battery to the persona axis
(full $2\times2\times2$ coverage): (i)~personas produce no spillover into the
judged Cultural Sensitivity dimension (pooled $d = -0.10$ to $-0.01$, all
domain cells n.s.; the personas were designed as a general cultural-perspective
injection, not as an instrument for this specific rubric item); (ii)~the
solo-persona novelty gain is genuine --- it survives length control
($+0.465$, $p < 0.001$) and is uniform across all six domains ($+0.36$ to
$+0.57$) against a uniform usefulness cost ($-0.96$ to $-1.10$); (iii)~teams
destroy that gain (structure $\times$ persona interactions: Novelty $-0.693$,
Usefulness $+0.711$, both $p < 0.001$, Holm-corrected; Holm, 1979);
(iv)~persona answers are judged \emph{more} consistently (panel ICC(3,9)
$0.933$ vs.\ $0.918$), so the effect is not a measurement artifact; and
(v)~personas roughly double outcome variance (solo SD $0.232 \to 0.397$),
consistent with a style-injection account. Length bias in LLM judging is a
documented threat~(Zheng et~al., 2023; Dubois et~al., 2024), which is why each
effect above was re-tested under length control before reporting.

\begin{quote}
\textbf{Strategic Implication (exploratory --- gate FAILED):} The contingency
picture above did \emph{not} survive the replication wave
(\cref{sec:wave}): on fresh problems the frontier advantage is uniform across
domains, the capability-gap correlation vanishes ($r = -0.21$ vs.\
$-0.978$), and the marketing variance advantage reverses. We retain it as a
hypothesis generator only; the decision framework
of~\cref{sec:decision-framework} rests on the replicated core claims, not on
these.
\end{quote}

\subsubsection{Out-of-sample confirmation: the replication wave (Protocol 12)}\label{sec:wave}

\textbf{Design.} Protocol~12 (pre-registered, OSF~\href{https://osf.io/jg49x}{jg49x},
before any wave generation) froze a binding upgrade rule: a Study-1
\emph{exploratory} finding becomes a headline claim only if its pre-registered
wave test passes at the pre-registered alpha and direction; the core claims are
re-tested with the same frozen instruments. The wave uses 30 fresh post-cutoff
business problems (a second $6 \times 5$ domain grid, zero overlap with
Study~1), the three core cells (C-S-N, C-T-N, F-S-N), and a family universe
doubled from 5 to 10 (adding Mistral, MiniMax, Gemma, GPT-OSS, and
ThinkingMachines): 2{,}400 generation runs, judged by the same frozen
three-judge panel under the same three verbatim prompt variants
(21{,}762 judgment cells; 2{,}418 analyzed runs). The analysis script was
frozen and hash-locked before judging; a disclosed numerical fallback (a
singular dual variance-component LMM, replaced under a pre-declared rule that
hard-requires bit-identical fixed effects) is recorded in the results artifacts
and our decision log (D52). Provider events are disclosed, not imputed: the
Gemini lane was deferred on a free-tier daily cap until a billing upgrade, and
the Kimi frontier lane was completed via a Together.ai routing deviation
(same pinned model version; D53) after Moonshot credit exhaustion.

\textbf{Gate outcomes.} \cref{tab:wave-gates} reports the complete-panel
results (the interim 9-family record of 2026-08-24 stands unchanged in the
wave report; the complete panel adds Gemini and the Kimi remainder).

\begin{table}[ht]
  \centering
  \caption{Confirmatory replication wave: gate outcomes on the complete
  10-family panel ($n = 2{,}418$ judged runs, $\alpha = 0.05$).}
  \label{tab:wave-gates}
  \begin{tabularx}{\textwidth}{@{}l X c@{}}
    \toprule
    \textbf{Gate} & \textbf{Result (complete panel)} & \textbf{Verdict} \\
    \midrule
    W1 team effect (C-S-N $\to$ C-T-N) & $+0.0168$, $p = 0.002$ & PASS \\
    W2 upgrade threshold (C-T-N $\to$ F-S-N) & $+0.0745$, $p < 0.0001$ & PASS \\
    W3 domain contingency & interaction LRT $p = 1.0$; directional signs $5/6$;
      \emph{no} positive domain (marketing flips negative) & FAIL \\
    W4 capability-gap mechanism & $r = -0.214$, $p = 0.55$ ($n = 10$ families;
      power for $\rho = 0.80$ is $0.83$) vs.\ Study-1 $r = -0.978$ & FAIL \\
    W5 marketing variance advantage & Brown--Forsythe $p = 0.795$,
      direction reversed (SD$_{\text{team}}$ $0.272$ $>$ SD$_{\text{solo}}$ $0.263$) & FAIL \\
    W6 economics & F-S-N $40.6$ $>$ C-T-N $21.2$ quality/\$-min; Pareto
      $=$ \{C-S-N, F-S-N\}; T-EGR $4.04$ (C-S-N), $0.96$ (C-T-N) & PASS \\
    W7 judge reliability (audit) & ICC(3,1) $0.471$--$0.622$ by domain & --- \\
    \bottomrule
  \end{tabularx}
\end{table}

\textbf{Reading.} The three core claims replicate out-of-sample on the full
family universe: (i)~cheap teams beat cheap solo by a small but real margin
(W1); (ii)~frontier solo beats cheap teams decisively on quality (W2, an
effect $\sim 4.4\times$ the team effect); (iii)~the cost--quality frontier is
bimodal --- cheap solo for volume economics (T-EGR $= 4.04$), frontier solo
for quality --- with cheap teams the dominated middle at current list prices
(W6; the ordering is unchanged under the MiniMax promotional-price sensitivity
vector, T-EGR $4.11$/$0.98$). All three exploratory upgrade candidates fail:
the ``marketing exception'' was a sample-specific fluctuation (marketing flips
negative, $\Delta = -0.075$), the near-perfect capability-gap correlation was
an $n = 5$ artifact, and teams do not reduce downside variance in any domain.
We read the failed gates as the design working as intended: the
pre-registration separated what survives contact with new data from what does
not~(Open Science Collaboration, 2015; Camerer et~al., 2018), and no post-hoc
relabeling is permitted in either direction.

\begin{quote}
\textbf{Key Takeaway:} the headline economics of this paper are
replication-grade; the contingency nuances are not. Practitioner rules should
be built on W1/W2/W6 only.
\end{quote}

\subsection{Study 2: Why Substitution Stops Short of the Frontier}\label{sec:study2}

Study~1 established that orchestrated lower-tier models approach but do not
reach the frontier (\cref{sec:frontier-gap}). Study~2 asks the follow-up a
budget owner actually faces: \emph{can a cheap refinement pass close that
residual gap --- and if it cannot, why not?} The answer is what separates a
risk-mitigation strategy from a replacement strategy.

\subsubsection{The refinement ceiling: closing the gap at low marginal cost}\label{sec:refinement}

A frozen, add-only refiner is applied to the CTP outputs. The pilot satisfied
the feasibility gate (G1, 100\% completion); the locked pool comprised all 171
matched CTP--FSN pairs (57 per family), of which the pilot covered 20 per
family and a pre-registered extension (registered before execution) covered
the remainder. The pre-registered superiority gate (G2: $d \ge 0.2$ with CI
excluding~0 in all three families, in the direction refined $>$ FSN) is
mechanically failed in all three families: kimi $d = -0.295$, claude $d =
-1.301$, qwen $d = -0.809$ (all CIs strictly negative after Holm). The
superiority claim is formally rejected.

Refinement nevertheless does substantial repair work. Measured as gap recovery
against each family's own pre-refinement CTP-vs-FSN gap, the pass recovers 60\%
of the gap for kimi, 8\% for claude (whose pre-refinement gap was already
small), and 64\% for qwen. The Novelty protection criterion (Protocol~09
criterion H2.2: no family loses more than $-0.10$ Novelty) is passed by all
three; Novelty actually \emph{increases} for qwen ($d = 1.10$,
$p_{\mathrm{FDR}} < 0.001$). Usefulness
and Elaboration show large gains in kimi and qwen ($d = 0.57$ to $1.35$, all
$p_{\mathrm{FDR}} < 0.001$) and are unchanged in claude --- which was already at
ceiling on those dimensions before refinement.

Length analysis confirms the dilution hypothesis is \emph{not} the driver: the
pooled Pearson correlation between word-count ratio (refined/original) and
composite-score change is $r = +0.67$, Spearman $\rho = +0.70$ ($p < 0.001$).
Longer refined answers score higher, not lower.

Cross-family comparisons (\cref{fig:cross-family}) remain negative after
refinement: no cross-family CTP pairing flips to superiority, confirming that
the within-family result is not an artifact of family selection.

\begin{quote}
\textbf{Strategic Implication:} A \$0.002 refinement pass delivers bounded
improvement --- it recovers most of the frontier gap on average but never
achieves paired superiority. For firms, this means refinement is a viable
risk-mitigation strategy (improve quality at lower cost) but not a replacement
strategy (it will not beat the frontier).
\end{quote}

\begin{figure}[ht]
  \centering
  \includegraphics[width=\textwidth]{fig4_gap_recovery.png}
  \caption{Study~2: Refinement Probe Results.}
  \label{fig:gap-recovery}
\end{figure}

\begin{figure}[ht]
  \centering
  \includegraphics[width=\textwidth]{fig7_cross_family_matrix.png}
  \caption{Cross-Family Comparison Matrix (Study~2).}
  \label{fig:cross-family}
\end{figure}

\subsubsection{Why orchestration buys coverage but not density}\label{sec:mechanism}

Thirty triples (10 per family) are coded under a 13-code codebook. Pooled
Cohen's $\kappa = 0.824$ (93.1\% agreement), well above the 0.70 gate. Six of
the thirteen codes fall below 0.70 (\texttt{diagnostic\_reframe} 0.54,
\texttt{named\_concepts} 0.49, \texttt{assumption\_test} 0.56,
\texttt{structured\_report} 0.59, \texttt{external\_authority} 0.56,
\texttt{redundancy\_bloat} 0.65), a base-rate artifact: percentage agreement on
these codes ranges from 82\% to 98\%, but $\kappa$ is depressed by near-zero
prevalence. Conclusions on below-threshold codes are disclosed as medium
reliability.

Of the six codes whose prevalence changed between original and refined, the
pattern differs by family. For kimi, all six changed codes are
\emph{structural scaffolding} codes: \texttt{operational\_parameters},
\texttt{structured\_report}, \texttt{phased\_timeline},
\texttt{tables\_matrices}, \texttt{quantified\_targets},
\texttt{granular\_ops\_detail}. For qwen, five of the six are structural
scaffolding (the same list minus \texttt{structured\_report}) and one ---
\texttt{diagnostic\_reframe} --- is an insight code. No other insight code
(\texttt{named\_concepts}, \texttt{assumption\_test},
\texttt{concise\_maxims}) moved in either family. The convergence metric compares, for each of the 13 codebook codes, how far
the original and refined lower-tier answers sit from the frontier answer's
profile: a code counts as ``toward'' if refinement reduces that distance and
``away'' if it increases it. Refinement does not consistently move lower-tier
outputs toward the frontier (kimi: toward~4 / away~6, net divergence; qwen:
5/4; claude: 0/2 with 11 codes flat --- claude's originals were already at the
codebook ceiling). Where outputs move away, refinement is adding structure the
frontier answers do not carry --- overshooting on coverage rather than gaining
insight density.

The residual frontier edge is precisely located. The only $\kappa$-reliable
frontier-vs-refined difference $\ge 0.30$ is \texttt{concise\_maxims}: qwen
frontier 1.00 vs.\ refined 0.55 (gap~0.45). The medium-reliability
\texttt{redundancy\_bloat} code is 0.00 in every frontier sample and 0.65--0.95
in every refined sample. Claude's original answers are already at the codebook
ceiling ($\ge 0.90$ on 11 of 13 codes); the number of changed codes between
original and refined is zero --- a saturation mechanism that explains the 8\%
gap recovery in~\cref{sec:refinement}.

\begin{quote}
\textbf{Strategic Implication:} Orchestration is a coverage play --- it adds
structure, length, and completeness. The frontier's advantage is a density
play --- compressed insight with zero waste. Add-only refinement cannot
manufacture density, which explains why lower-tier teams approach but do not surpass
frontier quality.
\end{quote}

\textbf{A note on replication scope.} Unlike Study~1 and Study~3, whose headline
claims were each stress-tested on fresh problems in a pre-registered replication
(\cref{sec:wave}, \cref{sec:study3conf}), Study~2 was not carried into a
confirmatory wave. Its two findings therefore rest on the original locked
sample: the refinement ceiling is a pre-registered gate result --- the G2
superiority test failed mechanically in all three families
(\cref{sec:refinement}) --- and the coverage-versus-density mechanism is a
mixed-methods coding result ($\kappa = 0.824$; \cref{sec:mechanism}). We report
Study~2 as a \emph{mechanistic explanation} of the Study~1 ceiling:
well-supported on its own sample, but not yet independently replicated. The
natural next step --- a \emph{subtractive} refiner permitted to cut as well as
add --- is named in \cref{sec:future}.

\begin{figure}[ht]
  \centering
  \includegraphics[width=\textwidth]{fig5_code_heatmap.png}
  \caption{Code Prevalence Heatmap: Original, Refined, and Frontier (Study~2).}
  \label{fig:heatmap}
\end{figure}

\subsection{Study 3: Who Should Design the Pipeline?}\label{sec:study3}

\textbf{Design.} Every pipeline in Studies 1--2 was \emph{human-designed}: the
builder $\to$ critic $\to$ refiner structure was our fixed research instrument
(the registered-but-unexecuted A7 arm isolated exactly this variable).
Study~3 (Protocol~13; operative pre-registration
OSF~\href{https://osf.io/bwc93}{bwc93} --- an initial v1 registration
was superseded by this instrument revision and withdrawn --- with
attrition clarification \href{https://osf.io/639eq}{639eq}; all
registrations timestamped before the corresponding computation)
executes A7 as a stand-alone study. Using
DeeWaanAI's production auto-designer, the model receives only the problem text and authors
a bespoke multi-step pipeline --- free choice of roles, step count, and
critic-gated if/while feedback loops --- which is then executed verbatim by
the cheap tier. Two arms per family: \textbf{D-SELF} (the cheap sibling
designs \emph{and} executes) and \textbf{D-FRONT} (the frontier flagship
designs; the cheap sibling executes). The 19 confirmatory problems
$\times$ 3 replications $\times$ 2 arms $\times$ 5 families $=$ 570 planned
runs; 544 completed, with 26 disclosed attrition events concentrated in the
Kimi lane (21 designer-stage failures --- the cheap designer returned empty or
truncated JSON on the production design prompt --- plus 4 execution failures
and 1 Claude designer failure). The pre-registered gates frame this flow:
feasibility gate G1 (per-family completion $\geq 80\%$ excluding
moderation-blocked tasks) passed for four families and failed strictly for
Kimi ($89/114 = 78\%$); per the amended G1 rule, the Kimi lane is retained
with the fully disclosed attrition above. The superiority-stop gate G2 (halt
only if D-FRONT trails the frozen C-T-N baseline by $d < -0.30$ on Novelty)
passed at pilot ($d = +0.42$), clearing the main run. Solutions were judged by the same frozen
three-judge panel under the three verbatim prompt variants (4{,}896 unique
judgments). The primary contrast (A) tests H-D1: D-FRONT vs.\ D-SELF on Novelty,
paired within family $\times$ problem, Holm-adjusted across families;
contrast B tests H-D2, both arms against the frozen human-designed C-T-N
baseline; an author-requested contrast informs H-D3, testing both arms
against the frozen frontier-solo (F-S-N) baseline, and a cost analysis
prices every run at the frozen price vector.

\textbf{Contrast A: authorship barely matters.}
\cref{tab:study3-authorship} reports the primary contrast. Frontier-designed
pipelines are statistically indistinguishable from self-designed ones in four
of five families; the lone significant effect favors frontier authorship for
Claude ($d = 0.78$, $p_{\text{Holm}} = 0.014$). Pooled across families, the
authorship gap on Novelty is $4.03$ vs.\ $3.99$ --- under $0.1$ scale points.
A post-hoc robustness check on the five-dimension composite (registered in
the Protocol-13 plan) leaves the conclusion unchanged: Claude is the only
family with a significant authorship effect on either DV (composite
$d = 0.35$, $p_{\text{Holm}} = 0.021$); every other family is null on both.

\begin{table}[ht]
  \centering
  \caption{Study 3 primary contrast: D-FRONT vs.\ D-SELF on Novelty, paired
  within family $\times$ problem ($n = 19$ pairs), Holm-adjusted across the
  five families.}
  \label{tab:study3-authorship}
  \begin{tabular}{@{}l c c c c c@{}}
    \toprule
    \textbf{Family} & $\Delta$ (D-FRONT $-$ D-SELF) & $t$ & $d$ &
      $p_{\text{raw}}$ & $p_{\text{Holm}}$ \\
    \midrule
    Claude & $+0.136$ & $3.00$ & $0.78$ & $0.003$ & $\mathbf{0.014}$ \\
    GPT-5.6 & $+0.049$ & $1.77$ & $0.22$ & $0.078$ & $0.310$ \\
    Qwen & $+0.021$ & $0.73$ & $0.10$ & $0.463$ & $1.0$ \\
    Gemini & $+0.021$ & $0.55$ & $0.08$ & $0.582$ & $1.0$ \\
    Kimi & $-0.038$ & $-0.64$ & $-0.12$ & $0.522$ & $1.0$ \\
    \bottomrule
  \end{tabular}
\end{table}

\textbf{Contrast B: both arms beat the human design.} Against the frozen
human-designed C-T-N pipeline, model-authored pipelines win in \emph{all five
families on both arms} (D-SELF $d = 0.15$--$0.67$; D-FRONT $d = 0.24$--$1.33$
on Novelty). The auto-designer also discovers structure we did not hand it:
98.9\% (D-SELF) and 100\% (D-FRONT) of designs include at least one
critic-gated if/while feedback loop, averaging $5.3$--$5.4$ designed steps and
$5.8$--$6.2$ executed steps after gate loops --- richer than our fixed
three-step instrument. Cheap designers iterate more at the design stage
(D-SELF consumed $34.8$M tokens vs.\ $24.3$M for D-FRONT across the study).

\textbf{Contrast vs.\ frontier solo: orchestration is family-specific, not
general.} Pooled, the arms are at parity with frontier solo (D-SELF
$d = -0.22$; D-FRONT $d = -0.05$ on Novelty) --- but the pooled figure hides
the story (\cref{tab:study3-econ}). Cheap auto-designed pipelines \emph{beat}
their own family's frontier flagship running solo for Kimi
($d = +0.85$/$+0.74$, $p_{\text{Holm}} = 0.006$/$0.013$), Qwen
($d = +0.55$/$+0.61$), and Gemini ($d = +0.46$/$+0.51$, n.s.\ after Holm);
they tie GPT-5.6-solo; and they lose catastrophically to Claude-opus solo
($d = -3.15$/$-2.84$). By domain, D-FRONT meets or beats frontier solo in
strategy, marketing, HR, and product, and loses in finance and operations ---
but this domain pattern does not survive within-family conditioning: the
moderator is the \emph{family} (how high the solo frontier sits), not the
task. Frontier authorship does not rescue the losing families: where the cheap
executor is the bottleneck, a frontier-designed plan cannot compensate.

\begin{table}[ht]
  \centering
  \caption{Study 3 economics: cost per completed solution (USD, frozen price
  vector v2026-08-23) and Novelty effect vs.\ frontier solo (F-S-N).
  $\Delta Q$ = Cohen's $d$ on Novelty vs.\ F-S-N. Positive $\Delta Q$ at
  cost ratio $< 1$ = strictly dominant orchestration.}
  \label{tab:study3-econ}
  \begin{tabular}{@{}l c c c c c@{}}
    \toprule
    \textbf{Family} & \textbf{F-S-N cost} & \multicolumn{2}{c}{\textbf{D-SELF}} &
      \multicolumn{2}{c}{\textbf{D-FRONT}} \\
    \cmidrule(lr){3-4}\cmidrule(lr){5-6}
    & & cost ($\times$F-S-N) & $\Delta Q$ & cost ($\times$F-S-N) & $\Delta Q$ \\
    \midrule
    Gemini & \$0.043 & \$0.010 ($0.24\times$) & $+0.46$ & \$0.025 ($0.59\times$) & $+0.51$ \\
    Qwen & \$0.020 & \$0.040 ($2.0\times$) & $+0.55$ & \$0.046 ($2.3\times$) & $+0.61$ \\
    Kimi & \$0.107 & \$0.377 ($3.5\times$) & $+0.85$ & \$0.274 ($2.6\times$) & $+0.74$ \\
    GPT-5.6 & \$0.180 & \$0.729 ($4.0\times$) & $-0.07$ & \$0.512 ($2.8\times$) & $+0.13$ \\
    Claude & \$0.266 & \$0.438 ($1.7\times$) & $-3.15$ & \$0.668 ($2.5\times$) & $-2.84$ \\
    \bottomrule
  \end{tabular}
\end{table}

\begin{figure}[ht]
  \centering
  \includegraphics[width=\textwidth]{}
  \caption{Study~3 family $\times$ domain landscape: the Novelty effect of
    cheap auto-designed pipelines versus frontier solo (F-S-N) across the five
    families and six business domains. The binding moderator is the
    \emph{family} --- how high its solo frontier ceiling sits --- not the task
    domain: the domain pattern does not survive within-family conditioning
    (\cref{sec:study3conf}).}
  \label{fig:study3-family-domain}
\end{figure}

\textbf{Decision rule.} Study~3 converts orchestration from a blanket
recommendation into a routing table. For Gemini, cheap auto-designed pipelines
are strictly dominant --- better \emph{and} 2--4$\times$ cheaper than frontier
solo. For Qwen and Kimi, orchestration buys the largest quality gains in the
study at 2--3.5$\times$ solo cost (which remains low in absolute terms).
For Claude, no cheap-executed pipeline approaches the flagship: buy solo.
For GPT-5.6, quality is a tie and solo is cheaper. And because authorship
effects are negligible (Contrast A), when orchestration is indicated the cheap
model can design for itself --- the frontier design fee buys nothing, except
where the cheap designer itself is fragile (21 of 57 Kimi self-design attempts
--- 37\% --- failed at the design stage, making frontier- or human-authored
templates the safer default for that family). \cref{sec:study3conf} stress-tests this table on new problems:
the extremes (Kimi, Claude) transport, the mid-tier rows do not, so the table
should be read as an illustration of the \emph{method} rather than a durable
lookup.

\begin{quote}
\textbf{Key Takeaway:} who \emph{designs} the pipeline does not matter ---
models author better pipelines than our human instrument, and cheap models
author as well as frontier ones. Whether to \emph{use} a pipeline depends on
the family: orchestration is a substitute for tier where the frontier ceiling
is low (Gemini, Qwen, Kimi) and a dominated cost where it is high (Claude).
\end{quote}

\subsubsection{Out-of-sample confirmation: does family moderation generalize? (Protocol 16)}\label{sec:study3conf}

\textbf{Design.} Study~3 estimated family-specific orchestration effects on
the 19 confirmatory problems. Protocol~16 (pre-registered
OSF~\href{https://osf.io/zkm8v}{zkm8v}, registered and timestamped before any
generation) tests whether those estimates transport to \emph{new} problems.
We ran the D-SELF arm only --- pre-declared, since Study~3 Contrast~A found
authorship effects negligible in 4/5 families --- on 30 fresh wave problems
$\times$ 3 replications $\times$ the same five core families $=$ 450 planned
runs; 441 completed with 9 disclosed attrition events (Gemini 90/90, Claude
89/90, Qwen 89/90, GPT-5.6 88/90, Kimi 85/90; feasibility gate G1 passed for
every family). The comparator is the \emph{frozen} wave frontier-solo (F-S-N)
baseline --- zero new frontier-solo generation, so the test is a pure
transport test of Study~3's effect sizes. Solutions were judged by the same
frozen three-judge panel under the three verbatim prompt variants
(3{,}963/3{,}969 valid unique judgments; 6 judge calls returned persistently
unparseable content on three Claude solutions and are disclosed as 0.15\%
judged-missing attrition). The analysis script was frozen and hash-locked
(SHA-256 recorded in the lab manifest) \emph{before} judging began. Five
binding gates translate the Study~3 point estimates into directional
predictions: R1 Kimi $d > 0$; R2 Qwen $d > 0$; R3 Gemini $d > 0$; R4 GPT-5.6
equivalence (TOST, 90\% CI on $d$ within $\pm 0.30$); R5 Claude $d < -1.0$.
The headline claim ``family moderation replicates'' requires
R1 $\wedge$ R2 $\wedge$ R3 $\wedge$ R5.

\textbf{Results: the extremes replicate; the middle does not.}
\cref{tab:study3conf-gates} reports the gate outcomes. Two of five gates
pass, and they are the two \emph{extreme} families: Kimi's cheap
auto-designed pipelines again beat frontier solo, by an even larger margin
than predicted ($d = +1.25$ vs.\ $+0.85$ predicted, one-sided
$p < 10^{-4}$), and Claude's again lose badly ($d = -1.90$ vs.\ $-3.15$
predicted, well past the $-1.0$ threshold). The three middle families fail:
Qwen's predicted gain vanishes ($d = -0.23$), and Gemini and GPT-5.6
\emph{reverse sign} --- their self-designed pipelines score significantly
\emph{below} frozen frontier solo on the new problems (Gemini $d = -0.83$;
GPT-5.6 $d = -0.88$, 90\% CI $[-1.32, -0.43]$, so the R4 equivalence gate
fails in the negative direction, not toward zero). The headline conjunction
is therefore \textbf{not replicated} (2/5 gates).

\begin{table}[ht]
  \centering
  \caption{Study 3 confirmatory replication (Protocol 16): binding gate
  outcomes. $d$ = Cohen's $d$ on Novelty, D-SELF vs.\ frozen wave F-S-N,
  problem-level Welch tests, $n = 30$ problem pairs per family. Predictions
  are the Study 3 point estimates.}
  \label{tab:study3conf-gates}
  \begin{tabular}{@{}l c c c c c@{}}
    \toprule
    \textbf{Gate} & \textbf{Family} & \textbf{Predicted} & \textbf{Observed $d$}
      & \textbf{90\% CI} & \textbf{Verdict} \\
    \midrule
    R1 & Kimi     & $d > 0$ ($+0.85$)          & $+1.25$ & $[+0.79, +1.72]$  & \textbf{PASS} \\
    R2 & Qwen     & $d > 0$ ($+0.55$)          & $-0.23$ & $[-0.66, +0.19]$  & fail \\
    R3 & Gemini   & $d > 0$ ($+0.46$)          & $-0.83$ & $[-1.28, -0.39]$  & fail (reversed) \\
    R4 & GPT-5.6  & $|d| < 0.30$               & $-0.88$ & $[-1.32, -0.43]$  & fail (negative) \\
    R5 & Claude   & $d < -1.0$ ($-3.15$)       & $-1.90$ & $[-2.41, -1.38]$  & \textbf{PASS} \\
    \bottomrule
  \end{tabular}
\end{table}

\textbf{Interpretation.} Family moderation is real but not a stable lookup
table. The qualitative ordering --- orchestration substitutes for tier where
the frontier ceiling is low and is a dominated cost where it is high ---
survives at its extremes (Kimi, Claude), where effects are large. But for
mid-tier families the sign and magnitude of the orchestration effect are
problem-set-dependent: point estimates from 19 problems did not transport to
30 new problems drawn from the same generator. Combined with fast model
turnover (the underlying model versions themselves change within months),
any static per-family routing table is obsolete on arrival. The correct
deliverable is therefore not a table but a \emph{method}: cheap-tier
benchmarking-as-routing --- test the user's actual model on the user's actual
problem type at trivial cost (a D-SELF probe run costs \$0.01--\$0.73 per
solution; \cref{tab:study3-econ}) and route on the result.

\textbf{Why we did not extend to new families.} A natural next step ---
repeating the study on the five additional families used in the wave
replication --- was considered and deliberately not registered. The
scientific question it would answer (does family moderation exist for new
families?) is already answered affirmatively by Studies 1--3; and
Protocol~16's central finding is that mid-tier effect sizes do not transport
even \emph{within} known families across problem sets. Five additional
families would therefore add point estimates subject to the same
non-transport, at roughly 450 runs and 4{,}000 judgments of cost, without
changing the methodological conclusion. We flag a bounded generalization
probe on new families as future work, positioned as a product-method
validation (benchmarking-as-routing) rather than a confirmatory experiment.

\begin{quote}
\textbf{Key Takeaway:} the Study 3 extremes replicate on new problems ---
cheap pipelines beat frontier solo where the frontier ceiling is low (Kimi,
$d = +1.25$) and lose where it is high (Claude, $d = -1.90$) --- but
mid-tier effects do not transport across problem sets. Orchestration routing
must be measured per deployment, not looked up: benchmark-as-routing, not a
routing table.
\end{quote}

\subsection{Summary of Key Findings}

\begin{table}[ht]
  \centering
  \footnotesize
  \caption{Summary of key findings across the three studies and their
  pre-registered replications.}
  \label{tab:summary-findings}
  \begin{tabular}{@{}>{\raggedright\arraybackslash}p{7.6cm} >{\raggedright\arraybackslash}p{4.6cm} >{\raggedright\arraybackslash}p{2.2cm}@{}}
    \toprule
    \textbf{Finding} & \textbf{Evidence} & \textbf{Section} \\
    \midrule
    Team structure helps lower-tier models more than frontier ($H_1$) & $d = -0.227$ & \cref{sec:interaction} \\
    Frontier $>$ cheaper overall ($H_2$) & $d = +0.294$ & \cref{sec:interaction} \\
    Teams slightly beat solos on average ($H_3$) & $d = +0.117$ & \cref{sec:interaction} \\
    Persona penalty on net ($H_4$, reversed) & $d = -0.691$ & \cref{sec:interaction} \\
    Structure $\times$ persona interaction ($H_5$, n.s.) & $F = 1.23$, $p = 0.27$ & \cref{sec:interaction} \\
    Persona-loaded cells Pareto-dominated & cost--quality plane & \cref{sec:persona-penalty} \\
    CTP $<$ FSN in every family & paired comparisons & \cref{sec:frontier-gap} \\
    Refinement recovers up to 64\% of gap (8\% where gap already small); G2 rejected & Protocol 09 tests (H2.1--H2.3) & \cref{sec:refinement} \\
    Orchestration buys coverage; frontier wins on density & $\kappa = 0.824$; 13 codes & \cref{sec:mechanism} \\
    Pareto frontier stable across $\phiBlend$ & \{C-S-P, C-S-N, C-T-N, F-S-N, F-T-N\} & \cref{sec:pareto} \\
    \emph{Replicated (wave):} team effect, upgrade threshold, bimodal economics & W1/W2/W6 PASS, 10 families, $n = 2{,}418$ & \cref{sec:wave} \\
    \emph{Exploratory (gate FAILED in wave):} domain contingency; capability-gap mechanism; marketing variance advantage & W3/W4/W5 FAIL & \cref{sec:exploratory-heterogeneity}, \cref{sec:wave} \\
    \emph{Exploratory (not wave-tested; wave excluded persona cells):} persona novelty gain is solo-only; teams squash it & D45 battery & \cref{sec:exploratory-heterogeneity} \\
    Design authorship negligible (4/5 families); auto-designed $>$ human-designed pipelines (5/5) & Study 3, Protocol 13 & \cref{sec:study3} \\
    Cheap auto-designed pipelines beat frontier solo for Gemini/Qwen/Kimi; lose for Claude; tie GPT & $d$ up to $+0.85$; cost 0.24--3.5$\times$ solo & \cref{sec:study3} \\
    \emph{Partially replicated:} extremes transport to new problems (Kimi $d = +1.25$; Claude $d = -1.90$); mid-tier families do not --- routing must be measured, not tabulated & Protocol 16, 2/5 gates, $n = 441$ & \cref{sec:study3conf} \\
    \bottomrule
  \end{tabular}
\end{table}

% ==================================================================
% 5  DISCUSSION
% ==================================================================
\section{Discussion}\label{sec:discussion}

\subsection{A Decision Framework for AI Investment}\label{sec:decision-framework}

The results support a contingency framework for AI investment rather than a
universal rule. For budget-constrained firms, \textbf{invest in the lower tier
first}: a single call from a cheaper model (C-S-N) achieves 98\% of frontier
quality at 23\% of the cost ($\TEGR$~6.15) --- the highest-return
configuration in the study. \textbf{Add the three-step pipeline only when
coverage is the target} (timelines, matrices, quantified plans): the cheaper
team (C-T-N) reaches 98.8\% of frontier quality at 88\% of the cost
($\TEGR$~1.49). \textbf{Do not
orchestrate at the frontier tier}: frontier team cells spend 3.5$\times$ the
frontier-solo cost for +0.024 points of quality ($\TEGR$~0.27), making them the
worst dollar in the experiment. \textbf{Use a low-cost refinement pass as a
risk-mitigation strategy} when absolute quality parity is the target: the
refined kimi point reaches parity at 78\% of F-S-N's cost, the lowest-cost route to
that milestone among the configurations tested --- noting that parity here is
on pooled means, while the paired within-family comparison remains negative
(\cref{sec:refinement}). \textbf{Avoid personas}: they degrade quality sharply
while reducing cost only marginally --- a value-destroying configuration.

\begin{figure}[ht]
  \centering
  \includegraphics[width=\textwidth]{fig20_contingency_rules.png}
  \caption{Practitioner contingency rules. Decision flowchart mapping budget
    tier and task type to the recommended orchestration cell, with expected
    T-EGR. Derived from the \cref{sec:decision-framework} contingency model.}
  \label{fig:contingency}
\end{figure}

\begin{table}[ht]
  \centering
  \small
  \caption{Managerial quick-reference: recommended configuration by task profile.}
  \label{tab:managerial}
  \begin{tabularx}{\textwidth}{@{}>{\raggedright\arraybackslash}p{3.4cm} l ccc >{\raggedright\arraybackslash}X@{}}
    \toprule
    \textbf{Task profile} & \textbf{Cell} & \textbf{Quality} & \textbf{Cost} & \textbf{T-EGR} & \textbf{Override when\ldots} \\
    \midrule
    Coverage-heavy (plans, timelines, matrices)
      & C-S-N & 98\% & 23\% & 6.15
      & absolute quality non-negotiable $\to$ F-S-N \\
    Density-critical (exec briefs, strategic principles)
      & F-S-N & 100\% & 100\% & 1.00
      & budget tight $\to$ C-S-N at 98\% \\
    Budget-minimum, quality-adequate
      & C-S-P & 92\% & 17\% & 8.60
      & novelty needed, substance traded \\
    Refinement as risk mitigation
      & Refined C-T-P & $\sim$100\% & 78\% & ---
      & never achieves paired superiority \\
    \textbf{Avoid}
      & F-T-N & 100.5\% & 353\% & 0.27
      & worst dollar in the experiment \\
    \textbf{Avoid}
      & F-T-P & 95.3\% & 365\% & 0.27
      & persona penalty on top of team overhead \\
    \bottomrule
  \end{tabularx}
\end{table}

Study~3 adds a family-conditional override to this framework
(\cref{sec:study3}, \cref{tab:study3-econ}): where the deployment stack is
Gemini, the auto-designed cheap pipeline is strictly dominant over frontier
solo (better novelty at $0.24\times$ the cost); for Qwen and Kimi the pipeline
is the quality-maximizing choice at 2--3.5$\times$ solo cost; for Claude the
solo flagship remains untouchable and pipelines are a dominated expense; and in
all cases the design fee is better spent on execution than on authorship ---
self-designed pipelines match frontier-designed ones in four of five families.

\subsection{The Density Ceiling: Why Additive Pipelines Hit a Wall}

The qualitative mechanism (\cref{sec:mechanism}) explains the pattern.
Orchestration is compositional --- add a critic, add a refiner, add a section,
add a table. Composition grows the answer but cannot compress it. The frontier's
residual advantage sits exactly where additive composition fails:
\texttt{concise\_maxims} (compressed, quotable strategic principles) and zero
\texttt{redundancy\_bloat}. These are the outputs of a single model that has
been selected or trained to produce dense prose. An add-only refiner, by design,
cannot delete the boilerplate the team structure produced. The falsifiable
prediction that follows is that a \emph{subtractive} refiner --- one permitted
to cut as well as add --- would narrow the density gap. This is the central
future-work experiment implied by our results.

\subsection{Contribution to the Business of AI Literature}

Our findings extend the author's prior lineage in two directions. The EGR metric
introduced in \citet{Hariri2026b} measured the cost performance of \emph{personas
in single models}; T-EGR extends it to teams and pipelines, preserving the ratio
form while re-basing on the frontier cell. The team-composition contingency
identified on verifiable tasks in \citet{Hariri2025b} reappears here on open-ended
tasks, but with the sign reversed for personas (harmful rather than neutral on
substance). The jagged-frontier finding of \citet{DellAcquaEtAl2023} is replicated
as a \emph{cell-level} result: lower-tier models inside the persona-free,
coverage-oriented region beat their cost-matched frontier alternatives; cheaper
models pushed outside by persona-loading or by density demands fall below the
frontier.

The test-time-compute prediction of \citet{SnellEtAl2024} is partially confirmed:
more inference compute (team, refinement) improves the cheaper model, but only up
to a parity ceiling, and only when priced in dollars does the tradeoff reveal its
practitioner shape. Our contribution is the dollarization of the compute-optimal
curve on open-ended tasks where no verifier exists.

\subsection{Methodological Lessons for AI Deployment Research}

Two lessons carry beyond this study. First, preregistered negative superiority
results can be the spine of a paper, not its appendix. The G2 mechanical failure
in~\cref{sec:refinement} is as informative as any positive finding, and
preregistration is what made it reportable as a primary result rather than as a
null. Second, token-invariant costing --- storing tokens and re-pricing through a
frozen vector --- is reusable infrastructure for the field: every reported dollar
is reproducible by any reader who supplies an alternative price vector
\citep{Diewert1976, SandveEtAl2013}.


% ==================================================================
% 6  MANAGERIAL IMPLICATIONS
% ==================================================================
\section{Managerial Implications}\label{sec:managerial}

This section translates the empirical findings into decisions a budget owner can
act on. Where \cref{sec:decision-framework} states the contingency framework as
a research result, the guidance below is written as practice: a sequenced budget
rule, an operating model, a risk register, and a sourcing decision. Every claim
carries its dollar figure first and its supporting statistic second, and points
to the quick-reference table (\cref{tab:managerial}) rather than repeating it.

\subsection{The Investment Ladder: A Sequenced Budget Rule}

The central managerial finding is that AI capability should be \emph{purchased in
a sequence}, not all at once. Spend down the following ladder and stop at the
rung where the next dollar buys less quality than it costs:

\begin{enumerate}[leftmargin=*,nosep]
  \item \textbf{Default to a single lower-tier call (C-S-N).} It delivers 98\%
        of frontier quality for 23\% of the price ($\TEGR$~6.15) --- the
        highest-return dollar in the study and the correct default for most
        open-ended knowledge work.
  \item \textbf{Add the three-step team only when coverage is the deliverable
        (C-T-N)} --- timelines, option matrices, quantified plans. It reaches
        98.8\% of frontier quality at 88\% of the cost ($\TEGR$~1.49): still
        positive-return, but the marginal dollar now buys \emph{coverage}, not
        quality.
  \item \textbf{Buy a refinement pass as cheap insurance (\$0.002).} It lifts a
        refined lower-tier answer to parity with the weakest frontier at 78\% of
        frontier-solo cost --- but it never achieves superiority. Treat it as a
        risk-mitigation rung, not a replacement rung.
  \item \textbf{Reserve frontier solo (F-S-N) for density-critical output} ---
        executive briefs and compressed strategic principles, where the residual
        frontier advantage actually lives (\cref{sec:mechanism}).
  \item \textbf{Never orchestrate at the frontier tier (F-T-N).} A frontier team
        costs 3.5$\times$ the frontier-solo price (353\% of cost) for negligible
        quality gain ($\TEGR$~0.27) --- the worst dollar in the experiment.
\end{enumerate}

\noindent\textbf{Strategic implication.} The budget question is not ``frontier or
cheap?'' but ``which rung, for which task?'' A firm that routes coverage work to
the lower tier and reserves the frontier for density captures most of the quality
at a fraction of the cost; a firm that pays frontier prices for team
orchestration overpays for output it could have bought cheaper.

\subsection{An Operating Model: Route by Measurement, Not by Table}

Study~3 changes how the ladder is operated. The economics of orchestration are
\emph{family-specific} (\cref{sec:study3}, \cref{tab:study3-econ}): the
auto-designed cheap pipeline strictly dominates frontier solo on a Gemini stack
at $0.24\times$ the cost, maximizes quality for Qwen and Kimi at $2$--$3.5\times$
solo cost, and is a dominated expense on Claude, where the solo flagship is
untouchable. The managerial consequence is direct: \textbf{there is no universal
routing table a firm can copy.} The operating model is a \emph{benchmark-and-route}
discipline --- on a small held-out set of the firm's own tasks, measure the
cheap-pipeline-versus-frontier-solo $\TEGR$ for the specific model family in the
stack, then route on the measured number.

Two further operating rules follow. First, \textbf{spend the design fee on
execution, not authorship}: self-designed pipelines matched frontier-designed
ones in four of five families, so a firm does not need to pay a frontier model to
author its workflows. Second, \textbf{treat routing as a recurring procurement
KPI}, re-measured whenever a model version or an API price changes --- because the
dollar figures that drive the ladder move with the market, even when the quality
ordering does not.

\subsection{Risk Register: What Can Go Wrong and How to Bound It}

\begin{itemize}[leftmargin=*,nosep]
  \item \textbf{Persona risk.} Role/persona prompting degrades substance while
        reducing cost only marginally (\cref{sec:persona-penalty}) --- a value-destroying
        configuration that is easy to adopt by accident because it ``feels'' more
        professional. \emph{Control:} audit persona usage and strip it wherever
        substance, not style, is graded.
  \item \textbf{Density risk.} Orchestration buys coverage but not density, and a
        refinement pass closes most of the gap without ever beating the frontier
        (\cref{sec:mechanism}, \cref{sec:refinement}). \emph{Control:} do not
        expect a cheap additive pipeline to substitute for compressed, quotable
        insight; buy the frontier for that class of output.
  \item \textbf{Replication risk --- honestly bounded.} The core economics
        replicated out-of-sample across 10 model families (Protocol~12), but the
        exploratory family-extremes did not transport cleanly (Protocol~16 passed
        2 of 5 pre-registered gates). \emph{Control:} treat the headline rules as
        bankable and the family-specific extremes as hypotheses to re-measure
        locally, not as guarantees. This is the payoff of pre-registration --- we
        can tell the manager which findings survived a fresh bet and which did not.
\end{itemize}

\subsection{Build-versus-Buy: A Capability-Sourcing Decision}

Frontier models are \emph{bought}: a fixed per-call cost for a given quality
level. Orchestration pipelines are \emph{built}: an upfront architectural
investment that lowers the marginal cost per call. The build path pays off when
(a)~the task is coverage-sensitive, (b)~the quality bar is ``good enough'' rather
than ``best available,'' and (c)~the firm can absorb the engineering cost of
maintaining a multi-step pipeline. When density is the binding constraint, the buy
path (frontier solo) remains optimal.

The decision statistic is $\TEGR$: \textbf{build when the measured $\TEGR$ on your
own workload exceeds 1, buy when the frontier's density premium is the thing you
actually need.} Critically, this is a \emph{per-task-family} decision, not a
firm-wide posture --- the same organization should build cheap pipelines for
coverage work and buy frontier calls for density work, and should expect the mix
to shift as model prices and capabilities move.

% ==================================================================
% 7  LIMITATIONS
% ==================================================================
\section{Limitations}\label{sec:limitations}

Six acknowledged limitations bound the claims. (1)~Evaluation is LLM-only; no
human panel was recruited. Our bias controls (blinding, order shuffle, disjoint
judge families, three prompt variants) follow the best current practice
\citep{ZhengEtAl2023, VergaEtAl2024} but do not substitute for human agreement
on a subset. (2)~The blended-$\phiBlend$ price assumption is a disclosed
limitation; we mitigate it with a pre-registered sensitivity panel and require
cross-$\phiBlend$ stability for any frontier-membership claim. (3)~The task
domain is English business strategy; generalization to code, scientific writing,
or non-English domains is an empirical question. (4)~Cross-family comparisons
(e.g.\ kimi vs.\ claude refinement rates) are descriptive and confounded with
family quality; only within-family comparisons support causal language. Two
further disclosures --- the small truncation audit and the team-cost retry gap
--- are stated but not corrected. On the latter: team intermediate outputs are
excluded from judged answers by design (only the final answer is scored) but
\emph{are} included in cost; the sole cost understatement comes from retry
runs whose ledger totals omit completed earlier steps (110 of 950 completed
team runs; the failed attempts themselves logged zero tokens and zero latency,
consuming nothing billable; median true-to-recorded cost ratio of 1.69 on
affected runs ---
roughly 6\% of team-cell cost on average). The gap mildly \emph{flatters} the
team cells, but it is immaterial against the cost differences that drive the
conclusions (teams cost $\sim$3.8$\times$ the solo call): no ordering can
flip.

A fifth limitation is practical: our costs are API-list prices, not
enterprise-negotiated rates. Firms with volume commitments or reserved-capacity
agreements may face different absolute prices, though the \emph{relative}
rankings (lower tier $\approx$ 23\% of frontier) are preserved as long as the
price ratio between tiers remains in the 3--8$\times$ range observed across the
five families at the time of the experiment. Similarly, pipeline latency --- the
wall-clock time of three sequential API calls versus one --- is not priced here;
latency-sensitive applications may prefer frontier solo even when the dollar
cost favors the lower-tier pipeline.

A sixth limitation concerns stimulus sampling breadth. The confirmatory
analysis rests on 19 business problems --- below the 30--40 stimulus range
recommended for maximal generalizability power \citep{WestfallKennyJudd2014}.
Although leave-one-problem-out cross-validation (19 folds)
confirmed sign stability of every reported effect (max $|\Delta d| = 0.02$),
and the within-problem replication is dense (100 runs/problem, 5
families), the 19 problems do not exhaust the universe of possible
business tasks. External validity --- the claim that orchestration
economics generalizes beyond these specific problems --- is therefore
bounded. Future work should replicate with larger, more diverse problem
corpora spanning additional domains, languages, and task complexities.

A seventh limitation concerns the replication wave's execution environment
(\cref{sec:wave}). Two provider-side events required disclosed handling: the
Gemini lane ran after a billing upgrade resolved a free-tier daily request
cap, and the Kimi frontier lane was completed via a Together.ai routing
deviation after Moonshot credit exhaustion (D53) --- the pinned model version
was unchanged, but within-family routing homogeneity for Kimi was broken
(cheap cells Moonshot-direct, frontier via Together; per-run route membership
is auditable in the ledger). Separately, the pre-registered dual
variance-component mixed model was numerically singular on the wave design
matrix; a pre-declared fallback rule (hard-requiring bit-identical fixed
effects) substituted a parsimonious random-effects structure (D52), so the
W3a likelihood-ratio statistic carries a REML-comparability caveat.

An eighth limitation is specific to Study~3 (\cref{sec:study3}). Attrition was concentrated in one family: 37\% of Kimi self-design
attempts (21 of 57) failed at the \emph{designer} stage (empty or
truncated structured output on the production auto-design prompt), and
total Kimi-lane attrition reached 22\% (25 of 114, including a provider
moderation block on one problem), so the Kimi cells rest on fewer
observations ($n = 89$ of 114) and the designer-fragility finding is
itself confounded with provider routing (Moonshot vs.\ Together; D58). Study~3
also instantiates design authorship through a \emph{single} production
auto-designer at its default decoding settings; other designers, temperatures,
or step caps might shift the authorship contrast, and our instrument
deliberately disallows human-in-the-loop approval steps, so the results speak
to fully automated pipelines only.

% ==================================================================
% 8  FUTURE DIRECTIONS
% ==================================================================
\section{Future Directions}\label{sec:future}

The first priority named in earlier drafts --- an out-of-sample confirmatory
replication wave --- is now complete (\cref{sec:wave}); it promoted the core
claims and rejected the exploratory upgrades. The second --- pipeline design
authorship --- is likewise now complete (Study~3, \cref{sec:study3}): it
executed the registered A7 arm and found authorship effects to be negligible,
with orchestration economics family-specific. Two candidates from the original roadmap are now closed by the results
themselves: \textbf{loop control} (Protocol~14) is superseded because
Study~3 found critic-gated loops in 98.9--100\% of auto-designed
pipelines yet outcomes remained family-determined, and \textbf{static
routing} (Protocol~15) is superseded because Protocol~16 showed mid-tier
effects do not transport across problem sets --- routing is a measured
per-deployment discipline (\cref{sec:study3conf}), not a learnable
lookup. Four experiments follow directly. First, a \textbf{bounded
generalization probe} on new model families, positioned as a validation
of the benchmarking-as-routing method rather than a confirmatory
experiment (\cref{sec:study3conf}). Second, a \textbf{subtractive
refiner} study that is permitted to cut as well as add: this isolates
the density mechanism identified in~\cref{sec:mechanism}. Third, a
\textbf{human-panel replication} on a subset of runs to bound LLM-judge
bias in our specific design. Fourth, a \textbf{second task domain} (code
generation, scientific writing) to test whether the
coverage-vs-density mechanism generalizes beyond business strategy.

% ==================================================================
% 9  CONCLUSION
% ==================================================================
\section{Conclusion}\label{sec:conclusion}

For organizations deploying LLMs for knowledge work, three investment rules
emerge from four pre-registered empirical bases: Study~1 (1{,}900 runs), its
Protocol~12 replication wave (2{,}418 analyzed runs), Study~3 (544 completed
runs), and the Protocol~16 replication (441 completed runs) --- together
spanning 10 model families and 30 fresh problems. \textbf{First}, at the lower-tier end,
model-tier substitution works: a single call from a lower-tier model
delivers 98\% of frontier quality at 23\% of the cost ($\TEGR$~6.15) --- the
highest-return configuration tested; the three-step team variant adds coverage
at 88\% of the cost ($\TEGR$~1.49). \textbf{Second}, at the frontier end, team
orchestration destroys value: it costs 3.5$\times$ the frontier-solo price for
negligible quality gain ($\TEGR$~0.27). \textbf{Third}, a low-cost refinement pass
(\$0.002) recovers most of the frontier gap but never beats it --- a
bounded-improvement result that makes it a viable risk-mitigation strategy, not
a replacement strategy. The residual frontier advantage is density (compressed
insight, zero waste), not coverage --- a capability gap that add-only pipelines
cannot bridge. Personas are value-destroying: they degrade quality sharply
($d = -0.691$) while reducing cost only marginally. The governing principle is contingency, not universality: match the
right pipeline architecture to the right model tier for the right task, and
invest in refinement before investing in a frontier team.

The broader implication is a build-versus-buy framework for AI capability.
Frontier models are \emph{bought} --- a fixed per-call cost for a given quality
level. Orchestration pipelines are \emph{built} --- an upfront architectural
investment that lowers marginal cost per call. Our results show that the build
path pays off when (a)~the task is coverage-sensitive, (b)~the quality bar is
``good enough'' rather than ``best available,'' and (c)~the firm can absorb the
engineering cost of maintaining a multi-step pipeline. When density is the
binding constraint --- when the output must compress insight into minimal prose
--- the buy path (frontier solo) remains optimal. The contingency is not a
limitation of the study; it is the finding.

For the field, the contribution is a dollarized, pre-registered, cost-matched
empirical test of the substitution claim that sits at the heart of the
multi-agent AI industry. The claim is partially true: structure does substitute
for tier, but only up to a ceiling set by density. We offer this ceiling, and
the contingency rules that surround it, as one empirical data point toward a
more rigorous economics of AI orchestration.

% ==================================================================
% APPENDIX A — EXECUTIVE SUMMARY (moved from front matter in v9)
% ==================================================================
\clearpage
\appendix

\section{Executive Summary for Practitioners}\label{sec:exec-summary}

\begin{mdframed}[style=execsummary]
\noindent\textbf{Decision context.} A firm that needs LLM-generated knowledge
work can either (a)~pay for a frontier model or (b)~use lower-tier models,
optionally orchestrated into a multi-step pipeline. Which dollar buys more
quality?

\medskip
\noindent\textbf{Bottom line.} Across 1{,}900 runs, five model families, and 19
business problems: a single call from a lower-tier model delivers
\textbf{98\% of frontier quality at 23\% of the cost} ($\TEGR$~6.15).
Orchestrating that model into a three-step builder--critic--refiner pipeline
adds coverage for about 3.8$\times$ that cost ($\TEGR$~1.49). Orchestrating a
frontier model costs \textbf{3.5$\times$ more for negligible gain}. A \$0.002
refinement pass recovers most of the remaining gap but never beats the
frontier.

\medskip
\noindent\textbf{Three rules for AI budget allocation:}
\begin{enumerate}[leftmargin=*,nosep]
  \item \textbf{cheaper tier, single call} $\to$ highest ROI
    ($\TEGR$~6.15). Add the three-step pipeline only when coverage is the
    target --- timelines, matrices, quantified plans --- at $\sim$3.8$\times$
    the solo cost ($\TEGR$~1.49).
  \item \textbf{Frontier tier (solo)} $\to$ best absolute quality per dollar
    when density matters (compressed insight, zero boilerplate).
  \item \textbf{Never add personas as a quality lever} --- they degrade judged
    quality sharply ($d = -0.691$) while reducing cost only marginally.
\end{enumerate}

\medskip
\noindent\textbf{Caveat.} Cheaper configurations buy \emph{coverage} (more
structure, more length); frontier models win on \emph{density} (insight per
paragraph). If the task rewards concise strategic principles, the frontier is
worth the premium. If the task rewards completeness, the lower-tier
configuration is the better bet.
\end{mdframed}

\textbf{From findings to product.} This research was conducted on, and now
feeds back into, the DeeWaanAI orchestration platform. Three findings have
already changed the product: (1)~the persona penalty ($d = -0.691$) removed
personas as a default quality lever; (2)~Study~3 established that pipeline
\emph{design authorship} is a negligible factor, so the platform's
auto-designer lets the cheap model design for itself rather than billing a
frontier design pass; and (3)~Protocol~16's partial replication converted the
static per-family routing table into \emph{benchmarking-as-routing} --- a
cheap per-deployment probe (test the user's model on the user's problem type,
then route), because mid-tier effects do not transport across problem sets.
The loop is deliberate: the platform supplies the research instrument, the
pre-registered findings supply the product's routing logic, and the
disclosure in \cref{sec:ai-disclosure} keeps the relationship transparent.

% ==================================================================
% APPENDIX B — PROTOCOL AND ARM REGISTRY
% ==================================================================
\clearpage
\section{Protocol and Arm Registry}\label{app:registry}

A \textbf{protocol} is a pre-registered experiment plan frozen before its data
are generated; an \textbf{arm} is one experimental condition within a protocol
(\cref{sec:method}). \cref{tab:protocol-registry} lists every numbered
protocol in the research program with its role and status;
\cref{tab:arm-registry} defines every arm label used in the text.
Protocols 05--08 were designed as \emph{modifications of the original
registration} rather than new studies: 05 corrected pilot issues before the
main run; 06 extended the interim-2 analysis with a pre-registered robustness
battery (items 6.1--6.7); 07 added post-lock exploratory statistics (items
7.1--7.5); 08 replaced the deductive codebook of Protocol 02 with an inductive
one. In the decision log and analysis artifacts these four appear as
``Amendments A1--A4''; the paper uses protocol numbers throughout.

\begin{table}[ht]
  \centering
  \caption{Protocol registry. ``Frozen'' = time-stamped before any data for
  that question were generated. OSF GUIDs are mapped in
  \cref{tab:osf-registrations}.}
  \label{tab:protocol-registry}
  \footnotesize
  \begin{tabular}{@{}l l l@{}}
    \toprule
    \textbf{Protocol} & \textbf{Role} & \textbf{Status} \\
    \midrule
    01 & Study 1 factorial pre-registration ($H_1$--$H_5$) & frozen; executed \\
    02 & Qualitative analysis (deductive codebook) & frozen; superseded by Protocol 08 \\
    03 & Sequential analysis (O'Brien--Fleming interims) & frozen; executed \\
    04 & Cost-effectiveness analysis (T-EGR, Pareto) & frozen; executed \\
    05 & Modification of Protocol 01: pilot corrections & frozen; executed \\
    06 & Modification of Protocol 01: extended interim-2 robustness battery (items 6.1--6.7) & frozen; executed \\
    07 & Modification of Protocol 01: exploratory statistics (items 7.1--7.5) & frozen; executed \\
    08 & Modification of Protocol 02: inductive qualitative coding & frozen; executed \\
    09 & Study 2: refinement probe & frozen; executed \\
    10 & Study 2: qualitative coding & frozen; executed \\
    11 & Study 2 extension: refinement (GPT, Gemini; cross-family refiner) & frozen; executed \\
    12 & Confirmatory replication wave (30 fresh problems, 10 families) & OSF jg49x; executed \\
    13 & Study 3: pipeline design authorship (D-SELF vs.\ D-FRONT) & OSF bwc93/639eq; executed \\
    14 & Loop control (critic-gated cycles) & drafted; superseded by Study 3 (D59) \\
    15 & Static routing from problem text & drafted; superseded by Protocol 16 (D59) \\
    16 & Study 3 confirmatory replication (family-moderation transport) & OSF zkm8v; executed \\
    \bottomrule
  \end{tabular}
\end{table}

\begin{table}[ht]
  \centering
  \caption{Arm registry. The eight factorial cells are tier (C = cheaper,
  F = frontier) $\times$ structure (S = solo, T = team) $\times$ persona
  (N = none, P = persona). A5--A8 are mechanism arms registered in Study 1;
  D-SELF/D-FRONT are the Study 3 arms.}
  \label{tab:arm-registry}
  \footnotesize
  \begin{tabular}{@{}l l l@{}}
    \toprule
    \textbf{Arm} & \textbf{Definition} & \textbf{Status} \\
    \midrule
    C-S-N & cheaper solo, no persona & executed (Study 1 + wave) \\
    C-S-P & cheaper solo, persona & executed (Study 1) \\
    C-T-N & cheaper team (builder $\to$ critic $\to$ refiner), no persona & executed (Study 1 + wave) \\
    C-T-P & cheaper team, persona & executed (Study 1) \\
    F-S-N & frontier solo, no persona (baseline) & executed (Study 1 + wave) \\
    F-S-P & frontier solo, persona & executed (Study 1) \\
    F-T-N & frontier team, no persona & executed (Study 1 + wave) \\
    F-T-P & frontier team, persona & executed (Study 1) \\
    \midrule
    A5 & critic-gated WHILE loop on the cheap team & registered; never executed in Study 1 \\
    A6 & cost-matched frontier-solo baseline & registered; never executed in Study 1 \\
    A7 & adaptive per-problem pipeline designer & registered; executed as Study 3 \\
    A8 & shuffled persona-to-role mapping & registered; never executed in Study 1 \\
    \midrule
    D-SELF & cheap sibling designs \emph{and} executes the pipeline & executed (Study 3 + Protocol 16) \\
    D-FRONT & frontier flagship designs; cheap sibling executes & executed (Study 3) \\
    \bottomrule
  \end{tabular}
\end{table}

% ==================================================================
% REFERENCES
% ==================================================================
\begin{thebibliography}{99}

\bibitem[Albers and Lakens(2018)]{AlbersLakens2018}
Albers, C., \& Lakens, D. (2018).
When power analyses based on pilot data are biased: Inaccurate effect size
  estimators and follow-up bias.
\textit{Journal of Experimental Social Psychology}, \textit{74}, 187--195.
\href{https://doi.org/10.1016/j.jesp.2017.09.004}{doi:10.1016/j.jesp.2017.09.004}

\bibitem[Borenstein et~al.(2009)]{BorensteinEtAl2009}
Borenstein, M., Hedges, L.~V., Higgins, J.~P.~T., \& Rothstein, H.~R. (2009).
\emph{Introduction to Meta-Analysis}. Wiley.
doi:10.1002/9780470743386.

\bibitem[Brown and Forsythe(1974)]{BrownForsythe1974}
Brown, M.~B., \& Forsythe, A.~B. (1974). Robust tests for the equality of
variances. \emph{Journal of the American Statistical Association}, 69(346),
364--367. doi:10.1080/01621459.1974.10482955.

\bibitem[Brynjolfsson et~al.(2023)]{BrynjolfssonLiRaymond2023}
Brynjolfsson, E., Li, D., \& Raymond, L.~R. (2023).
\textit{Generative AI at Work} (NBER Working Paper No.~31161).
\href{https://doi.org/10.3386/w31161}{doi:10.3386/w31161}

\bibitem[Camerer et~al.(2018)]{CamererEtAl2018}
Camerer, C.~F., Dreber, A., Holzmeister, F., Ho, T.-H., Huber, J.,
Johannesson, M., \ldots{} Wu, H. (2018).
Evaluating the replicability of social science experiments in
\textit{Nature} and \textit{Science} between 2010 and 2015.
\textit{Nature Human Behaviour}, 2(9), 637--644.
\href{https://doi.org/10.1038/s41562-018-0399-z}{doi:10.1038/s41562-018-0399-z}

\bibitem[Cohen(1960)]{Cohen1960}
Cohen, J. (1960).
A coefficient of agreement for nominal scales.
\textit{Educational and Psychological Measurement}, \textit{20}(1), 37--46.
\href{https://doi.org/10.1177/001316446002000104}{doi:10.1177/001316446002000104}

\bibitem[Dell'Acqua et~al.(2023)]{DellAcquaEtAl2023}
Dell'Acqua, F., McFowland~III, E., Mollick, E.~R., Lifshitz-Assaf, H.,
  Kellogg, K., Rajendran, S., Krayer, L., Candelon, F., \& Lakhani, K.~R.
  (2023).
\textit{Navigating the jagged technological frontier: Field experimental
  evidence of the effects of AI on knowledge worker productivity and quality}
  (HBS Working Paper No.~24-013).
\href{https://doi.org/10.2139/ssrn.4573321}{doi:10.2139/ssrn.4573321}

\bibitem[Dietterich(2000)]{Dietterich2000}
Dietterich, T.~G. (2000). Ensemble methods in machine learning. In
\emph{Multiple Classifier Systems (MCS 2000), Lecture Notes in Computer
Science, 1857}, 1--15. Springer. doi:10.1007/3-540-45014-9\_1.

\bibitem[Dubois et~al.(2024)]{DuboisEtAl2024}
Dubois, T., Liang, P., \& Hashimoto, T.~B. (2024). Length-controlled
AlpacaEval: A simple way to debias automatic evaluators. In
\emph{Proceedings of the 41st International Conference on Machine Learning
(ICML)}. arXiv:2404.04475.

\bibitem[Diewert(1976)]{Diewert1976}
Diewert, W.~E. (1976).
Exact and superlative index numbers.
\textit{Journal of Econometrics}, \textit{4}(2), 115--145.
\href{https://doi.org/10.1016/0304-4076(76)90009-9}{doi:10.1016/0304-4076(76)90009-9}

\bibitem[Hansen and Salamon(1990)]{HansenSalamon1990}
Hansen, L.~K., \& Salamon, P. (1990). Neural network ensembles.
\emph{IEEE Transactions on Pattern Analysis and Machine Intelligence},
12(10), 993--1001. doi:10.1109/34.58877.

\bibitem[Hariri(2025a)]{Hariri2025a}
Hariri, Y. (2025a).
\textit{LLM-Creativity} [dataset and paper]. Zenodo.
\href{https://doi.org/10.5281/zenodo.17407392}{doi:10.5281/zenodo.17407392}

\bibitem[Hariri(2025b)]{Hariri2025b}
Hariri, Y. (2025b).
\textit{LLM Team Composition Paradox} [dataset and paper]. Zenodo.
\href{https://doi.org/10.5281/zenodo.17490945}{doi:10.5281/zenodo.17490945}

\bibitem[Hariri(2026b)]{Hariri2026b}
Hariri, Y. (2026b).
\textit{Doing More with Less: How Cultural Personas Bridge the Cost--Performance
  Gap in LLMs}. Zenodo.
\href{https://doi.org/10.5281/zenodo.19533619}{doi:10.5281/zenodo.19533619}

\bibitem[Holm(1979)]{Holm1979}
Holm, S. (1979). A simple sequentially rejective multiple test procedure.
\emph{Scandinavian Journal of Statistics}, 6(2), 65--70.
doi:10.2307/4615733.

\bibitem[Kendall and Babington Smith(1939)]{KendallBabingtonSmith1939}
Kendall, M.~G., \& Babington Smith, B. (1939). The problem of $m$ rankings.
\emph{The Annals of Mathematical Statistics}, 10(3), 275--287.
doi:10.1214/aoms/1177732186.

\bibitem[Lakens(2013)]{Lakens2013}
Lakens, D. (2013).
Calculating and reporting effect sizes to facilitate cumulative science: A
  practical primer for $t$-tests and ANOVAs.
\textit{Frontiers in Psychology}, \textit{4}, 863.
\href{https://doi.org/10.3389/fpsyg.2013.00863}{doi:10.3389/fpsyg.2013.00863}

\bibitem[Liang et~al.(2024)]{LiangEtAl2024}
Liang, T., He, Z., Jiao, W., Wang, X., Wang, Y., Wang, R., Yang, Y., Shi, S.,
  \& Tu, Z. (2024).
Encouraging divergent thinking in large language models through multi-agent
  debate.
In \textit{Proceedings of the 2024 Conference on Empirical Methods in Natural
  Language Processing} (pp.~17889--17904).
\href{https://doi.org/10.18653/v1/2024.emnlp-main.992}{doi:10.18653/v1/2024.emnlp-main.992}

\bibitem[Open Science Collaboration(2015)]{OSC2015}
Open Science Collaboration. (2015).
Estimating the reproducibility of psychological science.
\textit{Science}, \textit{349}(6251), aac4716.
\href{https://doi.org/10.1126/science.aac4716}{doi:10.1126/science.aac4716}

\bibitem[Nosek et~al.(2018)]{NosekEtAl2018}
Nosek, B.~A., Ebersole, C.~R., DeHaven, A.~C., \& Mellor, D.~T. (2018).
The preregistration revolution.
\textit{PNAS}, \textit{115}(11), 2600--2606.
\href{https://doi.org/10.1073/pnas.1708274114}{doi:10.1073/pnas.1708274114}

\bibitem[Noy and Zhang(2023)]{NoyZhang2023}
Noy, S., \& Zhang, W. (2023).
Experimental evidence on the productivity effects of generative artificial
  intelligence.
\textit{Science}, \textit{381}(6654), 187--192.
\href{https://doi.org/10.1126/science.adh2586}{doi:10.1126/science.adh2586}

\bibitem[Sandve et~al.(2013)]{SandveEtAl2013}
Sandve, G.~K., Nekrutenko, A., Taylor, J., \& Hovig, E. (2013).
Ten simple rules for reproducible computational research.
\textit{PLoS Computational Biology}, \textit{9}(10), e1003285.
\href{https://doi.org/10.1371/journal.pcbi.1003285}{doi:10.1371/journal.pcbi.1003285}

\bibitem[Shinn et~al.(2023)]{ShinnEtAl2023}
Shinn, N., Cassano, F., Berman, E., Gopinath, A., Narasimhan, K., \& Yao, S.
  (2023).
Reflexion: Language agents with verbal reinforcement learning.
In \textit{Advances in Neural Information Processing Systems 36 (NeurIPS 2023)},
  pp.~8634--8652.

\bibitem[Shrout and Fleiss(1979)]{ShroutFleiss1979}
Shrout, P.~E., \& Fleiss, J.~L. (1979). Intraclass correlations: Uses in
assessing rater reliability. \emph{Psychological Bulletin}, 86(2), 420--428.
doi:10.1037/0033-2909.86.2.420.

\bibitem[Snell et~al.(2024)]{SnellEtAl2024}
Snell, C., Lee, J., Xu, K., \& Kumar, A. (2024).
Scaling LLM test-time compute optimally can be more effective than scaling model
  parameters.
arXiv:2408.03314.
\href{https://doi.org/10.48550/arXiv.2408.03314}{doi:10.48550/arXiv.2408.03314}

\bibitem[Steegen et~al.(2016)]{SteegenEtAl2016}
Steegen, S., Tuerlinckx, F., Gelman, A., \& Vanpaemel, W. (2016).
Increasing transparency through a multiverse analysis.
\textit{Perspectives on Psychological Science}, \textit{11}(5), 702--712.
\href{https://doi.org/10.1177/1745691616658637}{doi:10.1177/1745691616658637}

\bibitem[Tomlinson et~al.(2024)]{TomlinsonEtAl2024}
Tomlinson, B., Black, R.~W., Patterson, D.~J., \& Torrance, A.~W. (2024).
The carbon emissions of writing and illustrating are lower for AI than for
  humans.
\textit{Scientific Reports}, \textit{14}, 3732.
\href{https://doi.org/10.1038/s41598-024-54271-x}{doi:10.1038/s41598-024-54271-x}

\bibitem[Verga et~al.(2024)]{VergaEtAl2024}
Verga, P., Hofst\"atter, S., Althammer, S., Su, Y., Piktus, A.,
  Arkhangorodsky, A., Xu, M., White, N., \& Lewis, P. (2024).
Replacing judges with juries: Evaluating LLM generations with a panel of diverse
  models.
arXiv:2404.18796.
\href{https://doi.org/10.48550/arXiv.2404.18796}{doi:10.48550/arXiv.2404.18796}

\bibitem[Wang et~al.(2023)]{WangEtAl2023}
Wang, X., Wei, J., Schuurmans, D., Le, Q., Chi, E., Narang, S., Chowdhery, A.,
  \& Zhou, D. (2023).
Self-consistency improves chain of thought reasoning in language models.
In \textit{Proceedings of the Eleventh International Conference on Learning
  Representations} (ICLR 2023).

\bibitem[Wei et~al.(2022)]{WeiEtAl2022}
Wei, J., Wang, X., Schuurmans, D., Bosma, M., Ichter, B., Xia, F., Chi, E.,
  Le, Q., \& Zhou, D. (2022).
Chain-of-thought prompting elicits reasoning in large language models.
In \textit{Advances in Neural Information Processing Systems 35 (NeurIPS 2022)},
  pp.~24824--24837.

\bibitem[Westfall et~al.(2014)]{WestfallKennyJudd2014}
Westfall, J., Kenny, D.~A., \& Judd, C.~M. (2014).
Statistical power and optimal design in experiments in which samples of
  participants respond to samples of stimuli.
\textit{Journal of Experimental Psychology: General}, \textit{143}(5),
  2020--2045.
\href{https://doi.org/10.1037/xge0000014}{doi:10.1037/xge0000014}

\bibitem[Yao et~al.(2023)]{YaoEtAl2023}
Yao, S., Yu, D., Zhao, J., Shafran, I., Griffiths, T.~L., Cao, Y., \&
  Narasimhan, K. (2023).
Tree of thoughts: Deliberate problem solving with large language models.
In \textit{Advances in Neural Information Processing Systems 36 (NeurIPS 2023)}.

\bibitem[Zheng et~al.(2023)]{ZhengEtAl2023}
Zheng, L., Chiang, W.-L., Sheng, Y., Zhuang, S., Wu, Z., Zhuang, Y., Lin, Z.,
  Li, Z., Li, D., Xing, E.~P., Zhang, H., Gonzalez, J.~E., \& Stoica, I.
  (2023).
Judging LLM-as-a-judge with MT-Bench and Chatbot Arena.
In \textit{Advances in Neural Information Processing Systems 36 (NeurIPS 2023)},
  Datasets and Benchmarks Track.

\end{thebibliography}


\clearpage

\end{document}
